% !TEX root = ../Elements.tex
% Greek and English text last checked 16/03/2013
%%%%%%
% BOOK 4
%%%%%%
\pdfbookmark[0]{Book 4}{book4}
\pagestyle{empty}
\begin{center}
{\Huge ELEMENTS BOOK 4}\\
\spa\spa\spa
{\huge\it Construction of Rectilinear Figures In and}\\[0.5ex] {\huge\it Around Circles}
\end{center}
%\vspace{2cm}
%\epsfysize=5in
%\centerline{\epsffile{Book04/front.eps}}
\newpage

%%%%%%%
% Definitions
%%%%%%%
\pdfbookmark[1]{Definitions}{def4}
\pagestyle{fancy}
\cfoot{{\greekfont τ̀ηεπαγε}}
\lhead{\large{\greekfont ΣΤΟΙΘΕΙΩΝ γ̀γν{4}.}}
\rhead{\large ELEMENTS BOOK 4}
\begin{Parallel}{}{} 
\ParallelLText{
\begin{center}
\large{{\greekfont Ὅροι}.}
\end{center}\vspace*{-7pt}

\ggn{1}.~{\greekfont Σθῆμα εὐθύγραμμον εἰς σχῆμα εὐθύγραμμον ἐγγράφ-εσθαι λέγεται, ὅταν ἑκάστη τῶν τοῦ ἐγγραφομένου σθήματ-ος γωνιῶν
ἑκάστης πλευρᾶς τοῦ, εἰς ὃ ἐγγράφεται, ἅπτηται.}

\ggn{2}.~{\greekfont Σθῆμα δὲ ὁμοίως περὶ σχῆμα περιγράφεσθαι λέγεται, ὅταν ἑκάστη
πλευρὰ τοῦ περιγραφομένου ἑκάστης γωνίας τοῦ, περὶ ὃ περιγράφεται,
ἅπτηται.}

\ggn{3}.~{\greekfont Σθῆμα εὐθύγραμμον εἰς κύκλον ἐγγράφεσθαι λέγεται, ὅταν
ἑκάστη γωνία τοῦ ἐγγραφομένου ἅπτηται τῆς τοῦ κύκλου
περιφερείας.}

\ggn{4}.~{\greekfont Σθῆμα δὲ εὐθύγραμμον περὶ κύκλον περιγράφε-σθαι λέγεται,
ὅταν ἑκάστη πλευρὰ τοῦ περιγραφομένου ἐφάπτηται τῆς τοῦ κύκλου
περιφερείας.}

\ggn{5}.~{\greekfont Κύκλος δὲ εἰς σχῆμα ὁμοίως ἐγγράφεσθαι λέγεται, ὅταν
ἡ τοῦ κύκλου περιφέρεια ἑκάστης πλευρᾶς τοῦ, εἰς ὃ ἐγγράφεται,
ἅπτηται.}

\ggn{6}.~{\greekfont Κύκλος δὲ περὶ σχῆμα περιγράφεσθαι λέγεται, ὅταν ἡ τοῦ
κύκλου περιφέρεια ἑκάστης γωνίας τοῦ, περὶ ὃ περιγράφεται, ἅπτηται.}

\ggn{7}.~{\greekfont Εὐθεῖα εἰς κύκλον ἐναρμόζεσθαι λέγεται, ὅταν τὰ πέρατα
α ὐτῆς ἐπὶ τῆς περιφερείας ᾖ τοῦ κύκλου.}}

\ParallelRText{
\begin{center}
{\large Definitions}
\end{center}

1.~A rectilinear figure is said to be inscribed in a(nother) rectilinear figure
when  the respective angles of the inscribed figure touch the respective
sides of the (figure) in which it is inscribed.

2.~And, similarly, a  (rectilinear) figure is said to be circumscribed about a(nother rectilinear) figure when the respective sides of the circumscribed
(figure) touch the respective angles of the (figure) about which it is
circumscribed.

3.~A rectilinear figure is said to be inscribed in a circle when
each angle of the inscribed (figure) touches the circumference of the
circle.

4.~And a rectilinear figure is said to be circumscribed about a circle
when each side of the circumscribed (figure) touches the
circumference of the circle.

5.~And, similarly, a circle is said to be inscribed in a (rectilinear) figure
when the circumference of the circle touches each  side of
the (figure) in which it is inscribed.

6.~And a circle is said to be circumscribed about a rectilinear (figure)
when the circumference of the circle touches each  angle of the
(figure) about which it is circumscribed.

7.~A straight-line is said to be inserted into a circle when its extremities
are on the circumference of the circle.}
\end{Parallel}

%%%%%%
% Prop 4.1
%%%%%%
\pdfbookmark[1]{Proposition 4.1}{pdf4.1}
\begin{Parallel}{}{} 
\ParallelLText{
\begin{center}
{\large \ggn{1}.}
\end{center}\vspace*{-7pt}

{\greekfont Εἰς τὸν δοθέντα κύκλον τῇ δοθείσῃ εὐθείᾳ μὴ μείζονι οὔσῃ
τῆς τοῦ κύκλου διαμέτρου ἴσην εὐθεῖαν ἐναρμόσαι.}\\

\epsfysize=2.2in
\centerline{\epsffile{Book04/fig01g.eps}}

{\greekfont Ἔστω ὁ δοθεὶς κύκλος ὁ ΑΒΓ, ἡ δὲ δοθεῖσα εὐθεῖα μὴ
μείζων τῆς τοῦ κύκλου διαμέτρου ἡ Δ. δεῖ δὴ εἰς τὸν
ΑΒΓ κύκλον τῇ Δ εὐθείᾳ ἴσην εὐθεῖαν ἐναρμόσαι.}

{\greekfont Ἤθθω τοῦ ΑΒΓ κύκλου διάμετρος ἡ ΒΓ. εἰ μὲν
οὖν ἴση ἐστὶν ἡ ΒΓ τῇ Δ, γεγονὸς ἂν εἴη τὸ
ἐπιταθθέν; ἐνήρμοσται γὰρ εἰς τὸν ΑΒΓ κύκλον τῇ
Δ εὐθείᾳ ἴση ἡ ΒΓ. εἰ δὲ μείζων ἐστὶν ἡ ΒΓ τῆς
Δ, κείσθω τῇ Δ ἴση ἡ ΓΕ, καὶ κέντρῳ τῷ Γ διαστήματι
δὲ τῷ ΓΕ κύκλος γεγράφθω ὁ ΕΑΖ, καὶ ἐπεζεύθθω
ἡ ΓΑ.}

{\greekfont Ἐπεὶ οὖν το Γ σημεῖον κέντρον ἐστὶ τοῦ ΕΑΖ κύκλου, ἴση
 ἐστὶν ἡ ΓΑ τῇ ΓΕ. ἀλλὰ τῇ Δ ἡ ΓΕ ἐστιν ἴση;
καὶ ἡ Δ ἄρα τῇ ΓΑ ἐστιν ἴση.}

{\greekfont Εἰς ἄρα τὸν δοθέντα κύκλον τὸν ΑΒΓ τῇ δοθείσῃ εὐθείᾳ τῇ
Δ ἴση ἐνήρμοσται ἡ ΓΑ; ὅπερ ἔδει ποιῆσαι.}}

\ParallelRText{
\begin{center}
{\large Proposition 1}
\end{center}

To insert a straight-line equal to
a given straight-line into a circle, (the latter straight-line) not being greater than the diameter of the circle.

\epsfysize=2.2in
\centerline{\epsffile{Book04/fig01e.eps}}

Let $ABC$ be the given circle, and $D$ the given straight-line (which is) not greater than
the diameter of the circle. So it is required to insert a straight-line, equal to the
straight-line $D$, into the circle $ABC$.

Let a diameter $BC$ of circle $ABC$ be drawn.$^\dag$ Therefore, if
$BC$ is equal to $D$, (then) that (which)  was prescribed has taken place.
For the (straight-line) $BC$, equal to the straight-line $D$, has been inserted into the circle $ABC$.
And if $BC$ is greater than $D$ then let $CE$ be made equal to $D$ [Prop.~1.3], and let the circle $EAF$ be drawn with center $C$ and radius $CE$. And let $CA$ be joined.

Therefore, since the point $C$ is the center of  circle $EAF$, $CA$ is equal to $CE$. But,
$CE$ is equal to $D$. Thus, $D$ is also equal to $CA$.

Thus, $CA$, equal to the given straight-line $D$,  has been inserted into
the given circle $ABC$. (Which is) the very thing it was required to do.}
\end{Parallel}
{\footnotesize \noindent$^\dag$ Presumably, by finding the center of the circle [Prop.~3.1], and then
drawing a line through it.}

%%%%%%
% Prop 4.2
%%%%%%
\pdfbookmark[1]{Proposition 4.2}{pdf4.2}
\begin{Parallel}{}{} 
\ParallelLText{
\begin{center}
{\large \ggn{2}.}
\end{center}\vspace*{-7pt}

{\greekfont Εἰς τὸν δοθέντα κύκλον τῷ δοθέντι τριγώνῳ ἰσογώνιον τρίγωνον
ἐγγράψαι.}

\epsfysize=2in
\centerline{\epsffile{Book04/fig02g.eps}}

{\greekfont Ἔστω ὁ δοθεὶς κύκλος ὁ ΑΒΓ, τὸ δὲ δοθὲν τριγωνον τὸ ΔΕΖ;
δεῖ δὴ εἰς τὸν ΑΒΓ κύκλον τῷ ΔΕΖ τριγώνῳ ἰσογώνιον
τρίγωνον ἐγγράψαι.}

{\greekfont Ἤθθω τοῦ ΑΒΓ κύκλου ἐφαπτομένη ἡ ΗΘ κατὰ τὸ Α,
καὶ συνεστάτω πρὸς τῇ ΑΘ εὐθείᾳ καὶ τῷ πρὸς α ὐτῇ σημείῳ
τῷ Α τῇ ὑπὸ ΔΕΖ γωνίᾳ ἴση ἡ ὑπὸ
ΘΑΓ, πρὸς δὲ τῇ ΑΗ εὐθείᾳ καὶ τῷ πρὸς α ὐτῇ
σημείῳ τῷ Α τῇ ὑπὸ ΔΖΕ [γωνίᾳ] ἴση ἡ ὑπὸ ΗΑΒ, καὶ
ἐπεζεύθθω ἡ ΒΓ.}

{\greekfont Ἐπεὶ οὖν κύκλου τοῦ ΑΒΓ ἐφάπτεταί τις εὐθεῖα ἡ ΑΘ, καὶ
ἀπὸ τῆς κατὰ τὸ Α ἐπαφῆς εἰς τὸν κύκλον διῆκται εὐθεῖα
ἡ ΑΓ, ἡ ἄρα ὑπὸ ΘΑΓ ἴση ἐστὶ τῇ ἐν τῷ ἐναλλὰχ τοῦ
κύκλου τμήματι γωνίᾳ τῇ ὑπὸ ΑΒΓ. ἀλλ ἡ ὑπὸ ΘΑΓ τῇ
ὑπὸ ΔΕΖ ἐστιν ἴση; καὶ ἡ ὑπὸ ΑΒΓ ἄρα γωνία τῇ ὑπὸ
ΔΕΖ ἐστιν ἴση. διὰ τὰ α ὐτὰ δὴ καὶ ἡ ὑπὸ ΑΓΒ τῇ ὑπὸ
ΔΖΕ ἐστιν ἴση; καὶ λοιπὴ ἄρα ἡ ὑπὸ ΒΑΓ λοιπῇ τῇ ὑπὸ
ΕΔΖ ἐστιν ἴση [ἰσογώνιον ἄρα ἐστὶ τὸ ΑΒΓ τρίγωνον τῷ
ΔΕΖ τριγώνῳ, καὶ ἐγγέγραπται εἰς τὸν ΑΒΓ κύκλον].}

{\greekfont Εἰς τὸν δοθέντα ἄρα κύκλον τῷ δοθέντι τριγώνῳ
ἰσογώνιον τρίγωνον ἐγγέγραπται; ὅπερ ἔδει ποιῆσαι.}}

\ParallelRText{
\begin{center}
{\large Proposition 2}
\end{center}

To inscribe a triangle, equiangular with a given triangle, in a given circle.

\epsfysize=2in
\centerline{\epsffile{Book04/fig02e.eps}}

Let $ABC$ be the given circle, and $DEF$ the given triangle. So it is required
to inscribe a triangle, equiangular with triangle $DEF$, in circle $ABC$.

Let $GH$ be drawn touching circle $ABC$ at $A$.$^\dag$ And let  (angle) $HAC$, equal to angle $DEF$,  be constructed
on the straight-line $AH$ at the point $A$ on it, and (angle) $GAB$, equal to
[angle] $DFE$,  on the straight-line $AG$ at the point $A$ on it
 [Prop.~1.23]. And let $BC$ be joined.
 
 Therefore, since some straight-line $AH$ touches the circle $ABC$, and the
 straight-line $AC$ has been drawn across (the circle) from the point of
 contact $A$, (angle) $HAC$ is thus equal to the angle $ABC$ in the alternate
 segment of the circle [Prop.~3.32]. But, $HAC$ is equal to $DEF$.
 Thus, angle $ABC$ is also equal to $DEF$. So, for the same (reasons),
 $ACB$ is also equal to $DFE$.  Thus, the remaining (angle) $BAC$ is equal
 to the remaining (angle) $EDF$ [Prop.~1.32]. [Thus, triangle
 $ABC$ is equiangular with triangle $DEF$, and has been inscribed in circle
 $ABC$].
 
 Thus, a triangle, equiangular with the given triangle, has been inscribed in
 the given circle. (Which is) the very thing it was required to do.}
\end{Parallel}
{\footnotesize \noindent$^\dag$ See the footnote
to Prop.~3.34.}

%%%%%%
% Prop 4.3
%%%%%%
\pdfbookmark[1]{Proposition 4.3}{pdf4.3}
\begin{Parallel}{}{} 
\ParallelLText{
\begin{center}
{\large \ggn{3}.}
\end{center}\vspace*{-7pt}

{\greekfont Περὶ τὸν δοθέντα κύκλον τῷ δοθέντι τριγώνῳ ἰσογώνιον τρίγωνον περιγράψαι.}

\epsfysize=1.8in
\centerline{\epsffile{Book04/fig03g.eps}}

{\greekfont Ἔστω ὁ δοθεὶς κύκλος ὁ ΑΒΓ, τὸ δὲ δοθὲν τρίγωνον τὸ ΔΕΖ; δεῖ
δὴ περὶ τὸν ΑΒΓ κύκλον τῷ ΔΕΖ τριγώνῳ ἰσογώνιον τρίγωνον
περιγράψαι.}

{\greekfont Ἐκβεβλήσθω ἡ ΕΖ ἐφ ἑκάτερα τὰ μέρη κατὰ τὰ Η, Θ σημεῖα,
καὶ εἰλήφθω τοῦ ΑΒΓ κύκλου κέντρον τὸ Κ, καὶ διήθθω, ὡς
ἔτυθεν, εὐθεῖα ἡ ΚΒ, καὶ συνεστάτω πρὸς τῇ ΚΒ εὐθείᾳ
καὶ τῷ πρὸς α ὐτῇ σημείῳ τῷ Κ τῇ μὲν ὑπὸ ΔΕΗ γωνίᾳ
ἴση ἡ ὑπὸ ΒΚΑ, τῇ δὲ ὑπὸ ΔΖΘ ἴση ἡ ὑπὸ ΒΚΓ, καὶ διὰ
τῶν Α, Β, Γ σημείων ἤθθωσαν ἐφαπτόμεναι τοῦ ΑΒΓ κύκλου
αἱ ΛΑΜ, ΜΒΝ, ΝΓΛ.}

{\greekfont Καὶ ἐπεὶ ἐφάπτονται τοῦ ΑΒΓ κύκλου αἱ ΛΜ, ΜΝ,
ΝΛ κατὰ τὰ Α, Β, Γ σημεῖα, ἀπὸ δὲ τοῦ Κ κέντρου ἐπὶ
τὰ Α, Β, Γ σημεῖα
ἐπεζευγμέναι εἰσὶν  αἱ ΚΑ, ΚΒ, ΚΓ, ὀρθαὶ ἄρα εἰσὶν αἱ
πρὸς τοῖς Α, Β, Γ σημείοις γωνίαι. καὶ ἐπεὶ τοῦ ΑΜΒΚ τετραπλεύρου αἱ τέσσαρες γωνίαι τέτρασιν ὀρθαῖς ἴσαι εἰσίν, ἐπειδήπερ καὶ
εἰς δύο τρίγωνα διαιρεῖται τὸ ΑΜΒΚ, καί εἰσιν ὀρθαὶ αἱ ὑπὸ
ΚΑΜ, ΚΒΜ γωνίαι, λοιπαὶ ἄρα αἱ ὑπὸ ΑΚΒ, ΑΜΒ δυσὶν ὀρθαῖς
ἴσαι εἰσίν. εἰσὶ δὲ καὶ αἱ ὑπὸ ΔΕΗ, ΔΕΖ δυσὶν ὀρθαῖς ἴσαι;
αἱ ἄρα ὑπὸ ΑΚΒ, ΑΜΒ ταῖς ὑπὸ ΔΕΗ, ΔΕΖ ἴσαι εἰσίν,
ὧν ἡ ὑπὸ ΑΚΒ τῇ ὑπὸ ΔΕΗ ἐστιν ἴση; λοιπὴ ἄρα ἡ ὑπὸ
ΑΜΒ λοιπῇ τῇ ὑπὸ ΔΕΖ ἐστιν ἴση.
 ὁμοίως δὴ
δειθθήσεται, ὅτι καὶ ἡ ὑπὸ ΛΝΒ τῇ ὑπὸ ΔΖΕ ἐστιν ἴση;
καὶ λοιπὴ ἄρα ἡ ὑπὸ ΜΛΝ [λοιπῇ] τῇ ὑπὸ ΕΔΖ ἐστιν
ἴση. ἰσογώνιον ἄρα ἐστὶ τὸ ΛΜΝ τρίγωνον τῷ ΔΕΖ
τριγώνῳ; καὶ περιγέγραπται περὶ τὸν ΑΒΓ κύκλον.}

{\greekfont Περὶ τὸν δοθέντα ἄρα κύκλον τῷ δοθέντι τριγώνῳ ἰσογώνιον
τρίγωνον περιγέγραπται; ὅπερ ἔδει ποιῆσαι.}}

\ParallelRText{
\begin{center}
{\large Proposition 3}
\end{center}

To circumscribe a triangle, equiangular with a given triangle, about a given
circle.

\epsfysize=1.8in
\centerline{\epsffile{Book04/fig03e.eps}}

Let $ABC$ be the given circle, and $DEF$ the given triangle. So it is required
to circumscribe a triangle, equiangular with triangle  $DEF$, about circle $ABC$.

Let $EF$ be produced in each direction to points $G$ and 
$H$. And
let the center $K$ of circle $ABC$ be found [Prop.~3.1].
And let the straight-line $KB$ be drawn, at random,  across ($ABC$).
And let (angle) $BKA$, equal to angle $DEG$, be constructed on the straight-line $KB$ at the point $K$ on it, and (angle) $BKC$, equal to
$DFH$ [Prop.~1.23]. And let the (straight-lines) $LAM$, $MBN$, and $NCL$ be 
drawn through the points $A$, $B$, and $C$ (respectively), touching the circle $ABC$.$^\dag$

And since $LM$, $MN$, and $NL$ touch circle $ABC$ at points $A$, $B$, and $C$
(respectively), and $KA$, $KB$, and $KC$ are joined from the center $K$
to  points $A$, $B$, and $C$ (respectively), the angles at points
$A$, $B$, and $C$ are thus right-angles [Prop.~3.18]. And
since the (sum of the) four angles of quadrilateral $AMBK$ is equal to
four right-angles,
inasmuch as $AMBK$ (can) also (be) divided into two triangles [Prop.~1.32],
and angles $KAM$ and $KBM$ are (both) right-angles, the (sum of the) remaining (angles), $AKB$ and $AMB$, is thus equal to two right-angles. 
And $DEG$ and $DEF$ is also equal to two right-angles [Prop.~1.13].
Thus, $AKB$ and $AMB$ is equal to $DEG$ and $DEF$, of which $AKB$ is
equal to $DEG$. Thus, the remainder $AMB$ is equal to the remainder $DEF$.
So, similarly, it can be shown that $LNB$ is also equal to $DFE$. Thus, the
remaining (angle) $MLN$ is also equal to the [remaining] (angle) $EDF$ [Prop.~1.32]. Thus,
triangle $LMN$ is equiangular with triangle $DEF$. And it has been drawn around
circle $ABC$.

Thus,  a triangle, equiangular with  the given triangle, has been circumscribed about the given
circle. (Which is) the very thing it was required to do.}
\end{Parallel}
{\footnotesize \noindent$^\dag$ See the footnote
to Prop.~3.34.}

%%%%%%
% Prop 4.4
%%%%%%
\pdfbookmark[1]{Proposition 4.4}{pdf4.4}
\begin{Parallel}{}{} 
\ParallelLText{
\begin{center}
{\large \ggn{4}.}
\end{center}\vspace*{-7pt}

{\greekfont Εἰς τὸ δοθὲν  τρίγωνον κύκλον ἐγγράψαι.}

\epsfysize=2in
\centerline{\epsffile{Book04/fig04g.eps}}

{\greekfont Ἔστω τὸ δοθὲν τρίγωνον τὸ ΑΒΓ; δεῖ δὴ εἰς τὸ ΑΒΓ τρίγωνον κύκλον ἐγγράψαι.}

{\greekfont Τετμήσθωσαν αἱ ὑπὸ ΑΒΓ, ΑΓΒ γωνίαι δίθα ταῖς ΒΔ, ΓΔ εὐθείαις,
καὶ συμβαλλέτωσαν ἀλλήλαις κατὰ τὸ Δ σημεῖον, καὶ ἤθθωσαν ἀπὸ
τοῦ Δ ἐπὶ τὰς ΑΒ, ΒΓ, ΓΑ 	εὐθείας κάθετοι αἱ ΔΕ, ΔΖ, ΔΗ.}

{\greekfont Καὶ ἐπεὶ ἴση ἐστὶν ἡ ὑπὸ ΑΒΔ γωνία τῇ ὑπὸ ΓΒΔ, ἐστὶ
δὲ καὶ ὀρθὴ ἡ ὑπὸ ΒΕΔ ὀρθῇ τῇ ὑπὸ ΒΖΔ ἴση, δύο
δὴ τρίγωνά ἐστι τὰ ΕΒΔ, ΖΒΔ τὰς δύο γωνίας ταῖς δυσὶ γωνίαις
ἴσας ἔθοντα καὶ μίαν πλευρὰν μιᾷ πλευρᾷ ἴσην τὴν ὑποτείνουσαν
ὑπὸ μίαν
τῶν ἴσων γωνιῶν κοινὴν α ὐτῶν τὴν
ΒΔ; καὶ τὰς λοιπὰς ἄρα πλευρὰς ταῖς λοιπαῖς πλευραῖς ἴσας
ἕχουσιν; ἴση ἄρα ἡ ΔΕ τῇ ΔΖ. διὰ τὰ α ὐτὰ δὴ καὶ ἡ
ΔΗ τῇ ΔΖ ἐστιν ἴση. αἱ τρεῖς ἄρα εὐθεῖαι αἱ ΔΕ, ΔΖ, ΔΗ
ἴσαι ἀλλήλαις εἰσίν; ὁ ἄρα κέντρῷ τῷ Δ καὶ διαστήματι
ἑνὶ τῶν Ε, Ζ, Η κύκλος γραφόμενος ἥχει καὶ διὰ τῶν λοιπῶν
σημείων καὶ ἐφάψεται τῶν ΑΒ, ΒΓ, ΓΑ εὐθειῶν διὰ τὸ ὀρθὰς εἶναι τὰς πρὸς τοῖς Ε, Ζ, Η σημείοις γωνίας. εἰ γὰρ τεμεῖ
α ὐτάς, ἔσται ἡ τῇ διαμέτρῳ τοῦ κύκλου πρὸς ὀρθὰς
ἀπ ἄκρας ἀγομένη ἐντὸς πίπτουσα τοῦ κύκλου; ὅπερ
ἄτοπον ἐδείθθη; οὐκ ἄρα ὁ κέντρῳ τῷ Δ διαστήματι
δὲ ἑνὶ τῶν Ε, Ζ, Η γραφόμενος κύκλος τεμεῖ τὰς ΑΒ, ΒΓ, ΓΑ εὐθείας; ἐφάψεται ἄρα α ὐτῶν, καὶ ἔσται ὁ 
κύκλος ἐγγεγραμμένος εἰς τὸ ΑΒΓ τρίγωνον. 
ἐγγεγράφθω ὡς ὁ ΖΗΕ.}

{\greekfont Εἰς ἄρα τὸ δοθὲν τρίγωνον τὸ ΑΒΓ κύκλος
ἐγγέγρ-απται ὁ ΕΖΗ; ὅπερ ἔδει ποιῆσαι.}}

\ParallelRText{
\begin{center}
{\large Proposition 4}
\end{center}

To inscribe a circle in a given triangle.

\epsfysize=2in
\centerline{\epsffile{Book04/fig04e.eps}}

Let $ABC$ be the given triangle. So it is required to inscribe
a circle in triangle $ABC$.

Let the angles $ABC$ and $ACB$ be cut in half by the
straight-lines $BD$ and $CD$ (respectively) [Prop.~1.9], and let them meet one another
at point $D$, and let $DE$, $DF$, and $DG$ be drawn from
point $D$, perpendicular to the straight-lines $AB$, $BC$, and $CA$ (respectively) [Prop.~1.12].

And since angle $ABD$ is equal to $CBD$, and the right-angle
$BED$ is also equal to the right-angle $BFD$, $EBD$ and $FBD$ are thus two triangles
having two angles equal to two angles, and one side equal to one side---the
(one) subtending one of the equal angles (which is) common to the (triangles)---(namely), $BD$. Thus, they will also have the 
remaining sides equal to the (corresponding) remaining sides [Prop.~1.26]. Thus, $DE$ (is)
equal to $DF$. So, for the same (reasons), $DG$ is also equal
to $DF$. Thus, the three straight-lines $DE$, $DF$, and $DG$ are equal to one another.
Thus, the circle drawn with center $D$, and radius one of $E$, $F$, or $G$,$^\dag$ will also
go through the remaining points, and will  touch the straight-lines $AB$, $BC$, and $CA$, on account of the angles at $E$, $F$, and $G$ being right-angles. For if it cuts (one of) them then it will be a (straight-line)
drawn at right-angles to a diameter of the circle, from its extremity, falling inside
the circle. The very thing was shown (to be) absurd [Prop.~3.16].
Thus, the circle drawn with center $D$, and radius one of $E$, $F$, or $G$,
does not cut the straight-lines $AB$, $BC$, and $CA$. Thus, it will
touch them, and will be the circle inscribed  in triangle $ABC$.
Let it be (so) inscribed, like $FGE$ (in the figure).

Thus, the circle $EFG$ has been inscribed in the given triangle $ABC$.
(Which is) the very thing it was required to do.}
\end{Parallel}
{\footnotesize \noindent$^\dag$ Here, and in the following propositions, it
is understood that the radius is actually one of $DE$, $DF$, or $DG$.}

%%%%%%
% Prop 4.5
%%%%%%
\pdfbookmark[1]{Proposition 4.5}{pdf4.5}
\begin{Parallel}{}{} 
\ParallelLText{
\begin{center}
{\large \ggn{5}.}
\end{center}\vspace*{-7pt}

{\greekfont Περὶ τὸ δοθὲν τρίγωνον κύκλον περιγράψαι.}

\epsfysize=1.05in
\centerline{\epsffile{Book04/fig05g.eps}}

{\greekfont Ἔστω τὸ δοθὲν τρίγωνον τὸ ΑΒΓ; δεῖ δὲ περὶ τὸ δοθὲν τρίγωνον
τὸ ΑΒΓ κύκλον περιγράψαι.}

{\greekfont Τετμήσθωσαν αἱ ΑΒ, ΑΓ εὐθεῖαι δίθα κατὰ τὰ Δ, Ε σημεῖα, καὶ
ἀπὸ τῶν Δ, Ε σημείων ταῖς ΑΒ, ΑΓ πρὸς ὀρθὰς ἤθθωσαν
αἱ ΔΖ, ΕΖ; συμπεσοῦνται δὴ ἤτοι ἐντὸς τοῦ ΑΒΓ τριγώνου
ἢ ἐπὶ τῆς ΒΓ εὐθείας ἢ ἐκτὸς τῆς ΒΓ.}

{\greekfont Συμπιπτέτωσαν πρότερον ἐντὸς κατὰ τὸ Ζ, καὶ ἐπεζεύθθ-ωσαν
αἱ ΖΒ, ΖΓ, ΖΑ. καὶ ἐπεὶ ἴση ἐστὶν ἡ ΑΔ τῇ ΔΒ, κοινὴ
δὲ καὶ πρὸς ὀρθὰς ἡ ΔΖ, βάσις ἄρα ἡ ΑΖ βάσει τῇ ΖΒ
ἐστιν ἴση. ὁμοίως δὴ δείχομεν, ὅτι καὶ ἡ ΓΖ τῇ ΑΖ ἐστιν
ἴση; ὥστε καὶ ἡ ΖΒ τῇ ΖΓ ἐστιν ἴση; αἱ τρεῖς
ἄρα αἱ ΖΑ, ΖΒ, ΖΓ ἴσαι  ἀλλήλαις εἰσίν. ὁ ἄρα κέντρῳ
τῷ Ζ διαστήματι δὲ ἑνὶ τῶν Α, Β, Γ κύκλος γραφόμενος ἥχει
καὶ διὰ τῶν λοιπῶν σημείων, καὶ ἔσται περιγεγραμμένος
ὁ κύκλος περὶ τὸ ΑΒΓ τρίγωνον. περιγεγράφθω ὡς ὁ ΑΒΓ.}

{\greekfont Ἀλλὰ δὴ αἱ ΔΖ, ΕΖ συμπιπτέτωσαν ἐπὶ τῆς ΒΓ εὐθείας
κατὰ τὸ Ζ, ὡς ἔχει ἐπὶ τῆς δευτέρας καταγραφῆς, καὶ
ἐπεζεύθθω ἡ ΑΖ. ὁμοίως δὴ δείχομεν, ὅτι τὸ Ζ σημεῖον
κέντρον ἐστὶ τοῦ περὶ τὸ ΑΒΓ τρίγωνον περιγραφομένου κύκλου.}

{\greekfont Ἀλλὰ δὴ αἱ ΔΖ, ΕΖ συμπιπτέτωσαν ἐκτὸς τοῦ ΑΒΓ τριγώνου κατὰ
τὸ Ζ πάλιν, ὡς ἔχει ἐπὶ τῆς τρίτης καταγραφῆς, καί 
ἐπεζεύθθωσαν αἱ ΑΖ, ΒΖ, ΓΖ. καὶ ἐπεὶ πάλιν ἴση ἐστὶν ἡ
ΑΔ τῇ ΔΒ, κοινὴ δὲ καὶ πρὸς ὀρθὰς ἡ ΔΖ, βάσις ἄρα
ἡ ΑΖ βάσει τῇ ΒΖ ἐστιν ἴση. ὁμοίως δὴ δείχομεν, ὅτι
καὶ ἡ ΓΖ τῇ ΑΖ ἐστιν ἴση; ὥστε καὶ ἡ ΒΖ τῇ ΖΓ ἐστιν
ἴση; ὁ ἄρα [πάλιν] κέντρῳ τῷ Ζ διαστήματι δὲ ἑνὶ τῶν
ΖΑ, ΖΒ, ΖΓ κύκλος γραφόμενος ἥχει καὶ διὰ τῶν λοιπῶν σημείων,
καὶ ἔσται περιγεγραμμένος περὶ τὸ ΑΒΓ τρίγωνον.}

{\greekfont Περὶ τὸ δοθὲν ἄρα τρίγωνον κύκλος περιγέγραπται; ὅπερ ἔδει ποιῆσαι.}}

\ParallelRText{
\begin{center}
{\large Proposition 5}
\end{center}

To circumscribe a circle about a given  triangle.

\epsfysize=1.05in
\centerline{\epsffile{Book04/fig05e.eps}}

Let $ABC$ be the given triangle. So it is required to circumscribe
a circle about the given triangle $ABC$.

Let the straight-lines $AB$ and $AC$ be cut in half at points $D$ and $E$ (respectively) [Prop.~1.10]. And let $DF$ and $EF$ be drawn from 
points $D$ and $E$, at right-angles to $AB$ and $AC$ (respectively) [Prop.~1.11].
So  ($DF$ and $EF$) will surely either meet inside triangle $ABC$, on the straight-line
$BC$, or beyond $BC$.

Let them, first of all, meet inside (triangle $ABC$) at (point) $F$, and let $FB$, 
$FC$, and $FA$ be joined. And since $AD$ is equal to $DB$, and $DF$ is
common and at right-angles, the base $AF$ is thus equal to the base $FB$
[Prop.~1.4]. So, similarly, we can show that $CF$ is also equal to
$AF$. So that $FB$ is also equal to $FC$. Thus, the three (straight-lines)
$FA$, $FB$, and $FC$ are equal to one another. Thus, the circle drawn with
center $F$, and radius one of $A$, $B$, or $C$, will also go through the remaining points.
And the circle will be circumscribed about triangle $ABC$. Let it
be (so) circumscribed, like $ABC$ (in the first diagram from the left).

And so, let $DF$ and $EF$ meet on the straight-line $BC$ at (point) $F$, like
in the second diagram (from the left). And let $AF$ be joined.
So, similarly, we can show that point $F$ is the center of the circle circumscribed about triangle $ABC$.

And so, let $DF$ and $EF$ meet outside triangle $ABC$, again at (point) $F$, like
in the third diagram (from the left). And let $AF$, $BF$, and $CF$ be joined.
And, again, since $AD$ is equal to $DB$, and $DF$ is common and at right-angles,
the base $AF$ is thus equal to the base $BF$ [Prop.~1.4]. So, similarly,
we can show that $CF$ is also equal to $AF$. So that $BF$ is also equal to $FC$.
Thus, [again] the circle drawn with center $F$, and radius one of $FA$, $FB$, and
$FC$, will also go through the remaining points. And it will be circumscribed
about triangle $ABC$.

Thus, a circle has been circumscribed about the given triangle. (Which is)
the very thing it was required to do.}
\end{Parallel}

%%%%%%
% Prop 4.6
%%%%%%
\pdfbookmark[1]{Proposition 4.6}{pdf4.6}
\begin{Parallel}{}{} 
\ParallelLText{
\begin{center}
{\large \ggn{6}.}
\end{center}\vspace*{-7pt}

{\greekfont Εἰς τὸν δοθέντα κύκλον τετράγωνον ἐγγράψαι.}

{\greekfont Ἔστω ἡ δοθεὶς κύκλος ὁ ΑΒΓΔ; δεῖ δὴ εἰς τὸν ΑΒΓΔ κύκλον
τετράγωνον ἐγγράψαι.}

{\greekfont Ἤθθωσαν τοῦ ΑΒΓΔ κύκλου δύο διάμετροι πρὸς ὀρθὰς ἀλλήλαις
αἱ ΑΓ, ΒΔ, καὶ ἐπεζεύθθωσαν αἱ ΑΒ, ΒΓ, ΓΔ, ΔΑ.}\\

\epsfysize=2.2in
\centerline{\epsffile{Book04/fig06g.eps}}

{\greekfont Καὶ ἐπεὶ ἴση ἐστὶν ἡ ΒΕ τῇ ΕΔ; κέντρον γὰρ τὸ Ε; κοινὴ
δὲ καὶ πρὸς ὀρθὰς ἡ ΕΑ, βάσις ἄρα ἡ ΑΒ βάσει τῇ ΑΔ ἴση
ἐστίν. διὰ τὰ α ὐτὰ δὴ καὶ ἑκατέρα τῶν ΒΓ, ΓΔ ἑκατέρᾳ  τῶν ΑΒ, 
ΑΔ ἴση ἐστίν; ἰσόπλευρον ἄρα ἐστὶ τὸ ΑΒΓΔ τετράπλευρον. λέγω
δή, ὅτι καὶ ὀρθογώνιον. ἐπεὶ γὰρ ἡ ΒΔ εὐθεῖα  διάμετρός ἐστι
τοῦ ΑΒΓΔ κύκλου, ἡμικύκλιον ἄρα ἐστὶ τὸ ΒΑΔ; ὀρθὴ ἄρα
ἡ ὑπὸ ΒΑΔ γωνία. διὰ τὰ α ὐτὰ δὴ καὶ ἑκάστη τῶν ὑπὸ ΑΒΓ, ΒΓΔ, ΓΔΑ ὀρθή ἐστιν; ὀρθογώνιον ἄρα ἐστὶ τὸ ΑΒΓΔ τετράπλευρον.
ἐδείθθη δὲ καὶ ἰσόπλευρον; τετράγωνον ἄρα ἐστίν. καὶ
ἐγγέγραπται εἰς τὸν ΑΒΓΔ κύκλον.}

{\greekfont Εἰς ἄρα τὸν δοθέντα κύκλον τετράγωνον ἐγγέγραπται τὸ ΑΒΓΔ; ὅπερ
ἔδει ποιῆσαι.}}

\ParallelRText{
\begin{center}
{\large Proposition 6}
\end{center}

To inscribe a square in a given circle.

Let $ABCD$ be the given circle. So it is required to inscribe a square in circle
$ABCD$.

Let two diameters  of circle $ABCD$,  $AC$ and $BD$, be drawn at right-angles to one another.$^\dag$ And let $AB$, $BC$, $CD$, and $DA$ be joined.

\epsfysize=2.2in
\centerline{\epsffile{Book04/fig06e.eps}}

And since $BE$ is equal to $ED$, for $E$ (is) the center (of the circle), and $EA$
is common and at right-angles, the base $AB$ is thus equal to the base $AD$ [Prop.~1.4]. So, for the same (reasons), each of $BC$ and $CD$ is equal
to each of $AB$ and $AD$. Thus, the quadrilateral $ABCD$ is equilateral. So I say
that (it is) also right-angled. For since the straight-line $BD$ is a diameter
of circle $ABCD$, $BAD$ is thus a semi-circle. Thus, angle $BAD$ (is) a right-angle
[Prop.~3.31]. So, for the same (reasons), (angles)
$ABC$, $BCD$, and $CDA$ are also each right-angles. Thus, the quadrilateral $ABCD$
is right-angled. And it was also shown (to be) equilateral. Thus, it is a square [Def.~1.22]. And it has been inscribed in circle
$ABCD$.

Thus, the square $ABCD$ has been inscribed in the given circle. (Which is)
the very thing it was required to do.}
\end{Parallel}
{\footnotesize \noindent$^\dag$ Presumably, by finding the center of the circle [Prop.~3.1], 
drawing a line through it, and then drawing a second line through it, at right-angles to the first [Prop.~1.11].} 

%%%%%%
% Prop 4.7
%%%%%%
\pdfbookmark[1]{Proposition 4.7}{pdf4.7}
\begin{Parallel}{}{} 
\ParallelLText{
\begin{center}
{\large \ggn{7}.}
\end{center}\vspace*{-7pt}

{\greekfont Περὶ τὸν δοθέντα κύκλον τετράγωνον περιγράψαι.}

{\greekfont Ἔστω ὁ δοθεὶς κύκλος ὁ ΑΒΓΔ; δεῖ δὴ περὶ τὸν ΑΒΓΔ κύκλον
τετράγωνον περιγράψαι.}

{\greekfont Ἤθθωσαν τοῦ ΑΒΓΔ κύκλου δύο διάμετροι πρὸς ὀρθὰς ἀλλήλαις
αἱ ΑΓ, ΒΔ, καὶ διὰ τῶν Α, Β, Γ, Δ σημείων ἤθθωσαν ἐφαπτόμεναι
τοῦ ΑΒΓΔ κύκλου αἱ ΖΗ, ΗΘ, ΘΚ, ΚΖ.}

{\greekfont Ἐπεὶ οὖν ἐφάπτεται ἡ ΖΗ τοῦ ΑΒΓΔ κύκλου, ἀπὸ δὲ τοῦ Ε κέντρου ἐπὶ τὴν κατὰ τὸ Α ἐπαφὴν ἐπέζευκται ἡ ΕΑ, αἱ
ἄρα πρὸς τῷ Α γωνίαι ὀρθαί εἰσιν. διὰ τὰ α ὐτὰ δὴ καὶ αἱ
πρὸς τοῖς Β, Γ, Δ σημείοις γωνίαι ὀρθαί εἰσιν. καὶ ἐπεὶ
ὀρθή ἐστιν ἡ ὑπὸ ΑΕΒ γωνία, ἐστὶ δὲ ὀρθὴ καὶ ἡ ὑπὸ
ΕΒΗ, παράλληλος ἄρα ἐστὶν ἡ ΗΘ τῇ ΑΓ. διὰ τὰ α ὐτὰ
δὴ καὶ ἡ ΑΓ τῇ ΖΚ ἐστι παράλληλος. ὥστε καὶ ἡ ΗΘ τῇ
ΖΚ ἐστι παράλληλος. ὁμοίως δὴ δείχομεν, ὅτι καὶ ἑκατέρα
τῶν ΗΖ, ΘΚ τῇ ΒΕΔ ἐστι παράλληλος. παραλληλόγραμμα ἄρα ἐστὶ
τὰ ΗΚ, ΗΓ, ΑΚ, ΖΒ, ΒΚ; ἴση ἄρα ἐστὶν ἡ μὲν ΗΖ τῇ ΘΚ,
ἡ δὲ ΗΘ τῇ ΖΚ. καὶ ἐπεὶ ἴση ἐστὶν ἡ ΑΓ τῇ ΒΔ, ἀλλὰ
καὶ ἡ μὲν ΑΓ ἑκατέρᾳ τῶν ΗΘ, ΖΚ, ἡ δὲ ΒΔ ἑκατέρᾳ
τῶν ΗΖ, ΘΚ ἐστιν ἴση [καὶ ἑκατέρα ἄρα τῶν ΗΘ, ΖΚ ἑκατέρᾳ
τῶν ΗΖ, ΘΚ ἐστιν ἴση], ἰσόπλευρον ἄρα ἐστὶ τὸ ΖΗΘΚ
τετράπλευρον. λέγω δή, ὅτι
καὶ ὀρθογώνιον. ἐπεὶ γὰρ παραλληλόγραμμόν ἐστι τὸ ΗΒΕΑ,
καί ἐστιν ὀρθὴ ἡ ὑπὸ ΑΕΒ, ὀρθὴ ἄρα καὶ ἡ ὑπὸ ΑΗΒ. ὁμοίως
δὴ δείχομεν, ὅτι καὶ αἱ πρὸς τοῖς Θ, Κ, Ζ γωνίαι ὀρθαί
εἰσιν. ὀρθογώνιον ἄρα ἐστὶ τὸ ΖΗΘΚ. ἐδείθθη δὲ καὶ ἰσόπλευρον;
τετράγωνον ἄρα ἐστίν. καὶ περιγέγραπται περὶ τὸν ΑΒΓΔ κύκλον.}

\epsfysize=2.2in
\centerline{\epsffile{Book04/fig07g.eps}}

{\greekfont Περὶ τὸν δοθέντα ἄρα κύκλον τετράγωνον περιγέγραπται;
ὅπερ ἔδει ποιῆσαι.}}

\ParallelRText{
\begin{center}
{\large Proposition 7}
\end{center}

To circumscribe a square about a given circle.

Let $ABCD$ be the given circle. So it is required to circumscribe a square
about circle $ABCD$.
Let two diameters of circle $ABCD$,  $AC$ and $BD$,  be drawn at right-angles to one another.$^\dag$
And let $FG$, $GH$, $HK$, and $KF$ be drawn through points $A$, $B$, $C$, and
$D$ (respectively), touching circle $ABCD$.$^\ddag$

Therefore, since $FG$ touches circle $ABCD$, and $EA$ has been joined from the
center $E$ to the point of contact $A$, the angles at $A$ are thus  right-angles  
[Prop.~3.18]. So, for the same (reasons), the angles at points $B$, $C$, and $D$ are also
right-angles. And since angle $AEB$ is a right-angle, and $EBG$ is also a right-angle, $GH$ is thus parallel to $AC$ [Prop.~1.29]. So, for the same
(reasons), $AC$ is also parallel to $FK$. So that $GH$ is also parallel to
$FK$ [Prop.~1.30]. So, similarly, we can show that $GF$ and $HK$ are each parallel to
$BED$. Thus, $GK$, $GC$, $AK$, $FB$, and $BK$ are (all) parallelograms. Thus, 
$GF$ is equal to $HK$, and $GH$ to $FK$ [Prop.~1.34]. And since
$AC$ is equal to $BD$, but $AC$ (is) also (equal) to each of $GH$ and $FK$,
and $BD$ is equal to each of $GF$ and $HK$ [Prop.~1.34] [and each of $GH$ and $FK$ is thus equal to each of $GF$ and $HK$], the quadrilateral $FGHK$ is
thus equilateral.  So I say that (it is) also right-angled. For since $GBEA$
is a parallelogram, and $AEB$ is a right-angle, $AGB$ is thus also a right-angle [Prop.~1.34]. So, similarly, we can show that the angles at $H$, $K$, and
$F$ are also right-angles. Thus, $FGHK$ is right-angled. And it was also shown
(to be) equilateral. Thus,  it is a square [Def.~1.22].
And it has been circumscribed about circle $ABCD$.

\epsfysize=2.2in
\centerline{\epsffile{Book04/fig07e.eps}}

Thus, a square has been circumscribed about the given circle. (Which is)
the very thing it was required to do.}
\end{Parallel}
{\footnotesize \noindent$^\dag$ See the footnote to the previous
 proposition.\\
 $^\ddag$ See the footnote
to Prop.~3.34.}

%%%%%%
% Prop 4.8
%%%%%%
\pdfbookmark[1]{Proposition 4.8}{pdf4.8}
\begin{Parallel}{}{} 
\ParallelLText{
\begin{center}
{\large \ggn{8}.}
\end{center}\vspace*{-7pt}

{\greekfont Εἰς τὸ δοθὲν τετράγωνον κύκλον ἐγγράψαι.}

{\greekfont Ἔστω τὸ δοθὲν τετράγωνον τὸ ΑΒΓΔ; δεῖ δὴ εἰς τὸ ΑΒΓΔ
τετράγωνον κύκλον ἐγγράψαι.}

\epsfysize=2.2in
\centerline{\epsffile{Book04/fig08g.eps}}

{\greekfont Τετμήσθω ἑκατέρα τῶν ΑΔ, ΑΒ δίθα κατὰ τὰ Ε, Ζ σημεῖα, καὶ
διὰ μὲν τοῦ Ε ὁποτέρᾳ τῶν ΑΒ, ΓΔ παράλληλος ἤθθω
ὁ ΕΘ, διὰ δὲ τοῦ Ζ ὁποτέρᾳ τῶν ΑΔ, ΒΓ παράλληλος ἤθθω
ἡ ΖΚ; παραλληλόγραμμον ἄρα ἐστὶν ἕκαστον τῶν ΑΚ, ΚΒ,
ΑΘ, ΘΔ, ΑΗ, ΗΓ, ΒΗ, ΗΔ, καὶ αἱ ἀπεναντίον α ὐτῶν πλευραὶ
δηλονότι ἴσαι [εἰσίν]. καὶ ἐπεὶ ἴση ἐστὶν ἡ ΑΔ τῇ ΑΒ,
καί ἐστι τῆς μὲν ΑΔ ἡμίσεια ἡ ΑΕ, τῆς δὲ ΑΒ ἡμίσεια ἡ ΑΖ,
ἴση ἄρα καὶ ἡ ΑΕ τῇ ΑΖ; ὥστε καὶ αἱ ἀπεναντίον; ἴση ἄρα
καὶ ἡ ΖΗ τῇ ΗΕ. ὁμοίως δὴ δείχομεν, ὅτι καὶ ἑκατέρα
τῶν ΗΘ, ΗΚ
ἑκατέρᾳ τῶν ΖΗ, ΗΕ ἐστιν ἴση; αἱ τέσσαρες ἄρα αἱ ΗΕ, ΗΖ, ΗΘ, ΗΚ ἴσαι ἀλλήλαις [εἰσίν]. ὁ ἄρα κέντρῳ μὲν τῷ Η διαστήματι
δὲ ἑνὶ τῶν Ε, Ζ, Θ, Κ κύκλος γραφόμενος ἥχει καὶ διὰ τῶν λοιπῶν
σημείων; καὶ ἐφάψεται τῶν ΑΒ, ΒΓ, ΓΔ, ΔΑ εὐθειῶν διὰ
τὸ ὀρθὰς εἶναι τὰς πρὸς τοῖς Ε, Ζ, Θ, Κ γωνίας; εἰ γὰρ τεμεῖ
ὁ κύκλος τὰς ΑΒ, ΒΓ, ΓΔ, ΔΑ, ἡ τῇ διαμέτρῳ τοῦ κύκλου
πρὸς ὀρθὰς ἀπ ἄκρας ἀγομένη ἐντὸς πεσεῖται τοῦ κύκλου;
ὅπερ ἄτοπον ἐδείθθη. οὐκ
ἄρα ὁ κέντρῳ τῷ Η διαστήματι δὲ ἑνὶ τῶν Ε, Ζ, Θ, Κ
κύκλος γραφόμενος τεμεῖ τὰς ΑΒ, ΒΓ, ΓΔ, ΔΑ εὐθείας. ἐφάψεται
ἄρα α ὐτῶν καὶ ἔσται ἐγγεγραμμένος εἰς τὸ ΑΒΓΔ τετράγωνον.}

{\greekfont Εἰς  ἄρα τὸ δοθὲν τετράγωνον κύκλος ἐγγέγραπται; ὅπερ
ἔδει ποιῆσαι.}}

\ParallelRText{
\begin{center}
{\large Proposition 8}
\end{center}

To inscribe a circle in a given square.

Let the given square be $ABCD$. So it is required to inscribe a circle in
square $ABCD$.

\epsfysize=2.2in
\centerline{\epsffile{Book04/fig08e.eps}}

Let  $AD$ and $AB$ each be cut in half at points $E$ and $F$ (respectively)
[Prop.~1.10]. And let $EH$ be drawn through $E$, parallel
to either of $AB$ or $CD$, and let $FK$ be drawn through $F$, parallel to either of $AD$ or $BC$ [Prop.~1.31].
Thus, $AK$, $KB$, $AH$, $HD$, $AG$, $GC$, $BG$, and $GD$ are each parallelograms,
and their opposite sides [are] manifestly equal [Prop.~1.34].
And since $AD$ is equal to $AB$, and $AE$ is half of $AD$, and $AF$ half of
$AB$, $AE$ (is) thus also equal to $AF$. So that the opposite (sides are) also (equal).
Thus, $FG$ (is) also equal to $GE$. So, similarly, we can also show that each of
$GH$ and $GK$ is equal to each of $FG$ and $GE$. Thus, the four (straight-lines)
$GE$, $GF$, $GH$, and $GK$ [are] equal to one another. Thus, the circle drawn
with center $G$, and radius one of $E$, $F$, $H$, or $K$, will also go through the remaining points. And it will touch the straight-lines $AB$, $BC$, $CD$, and $DA$,
on account of the angles at $E$, $F$, $H$, and $K$ being right-angles. For if the
circle cuts $AB$, $BC$, $CD$, or $DA$, (then) a (straight-line) drawn at right-angles to a diameter of the circle, from its extremity, will fall inside the circle. 
The very thing was shown (to be) absurd [Prop.~3.16].  Thus, the
circle drawn with center $G$, and radius one of $E$, $F$, $H$, or $K$, does not
cut the straight-lines $AB$, $BC$, $CD$, or $DA$. Thus, it will touch them, and will
be inscribed in the square $ABCD$.

Thus, a circle has been inscribed in the given square. (Which is) the very thing it was required to do.}
\end{Parallel}

%%%%%%
% Prop 4.9
%%%%%%
\pdfbookmark[1]{Proposition 4.9}{pdf4.9}
\begin{Parallel}{}{} 
\ParallelLText{
\begin{center}
{\large \ggn{9}.}
\end{center}\vspace*{-7pt}

{\greekfont Περὶ τὸ δοθὲν τετράγωνον κύκλον περιγράψαι.}

{\greekfont Ἔστω τὸ δοθὲν τετράγωνον τὸ ΑΒΓΔ; δεῖ δὴ περὶ τὸ ΑΒΓΔ τετράγωνον κύκλον περιγράψαι.}

{\greekfont Ἐπιζευθθεῖσαι γὰρ αἱ ΑΓ, ΒΔ τεμνέτωσαν ἀλλήλας κατὰ τὸ Ε.}

\epsfysize=2.2in
\centerline{\epsffile{Book04/fig06g.eps}}

{\greekfont Καὶ ἐπεὶ ἴση ἐστὶν ἡ ΔΑ τῇ ΑΒ, κοινὴ δὲ ἡ ΑΓ, δύο
δὴ αἱ ΔΑ, ΑΓ δυσὶ ταῖς ΒΑ, ΑΓ ἴσαι εἰσίν; καὶ βάσις ἡ ΔΓ βάσει
τῇ ΒΓ ἴση; γωνία ἄρα ἡ ὑπὸ ΔΑΓ γωνίᾳ τῇ ὑπὸ ΒΑΓ ἴση
ἐστίν; ἡ ἄρα ὑπὸ ΔΑΒ γωνία δίθα τέτμ ηται ὑπὸ τῆς ΑΓ. ὁμοίως
δὴ δείχομεν, ὅτι καὶ ἑκάστη τῶν ὑπὸ ΑΒΓ, ΒΓΔ, ΓΔΑ δίθα
τέτμ ηται ὑπὸ τῶν ΑΓ, ΔΒ εὐθειῶν. καὶ ἐπεὶ ἴση
ἐστὶν ἡ ὑπὸ ΔΑΒ γωνία τῇ ὑπὸ ΑΒΓ, καί ἐστι τῆς μὲν ὑπὸ
ΔΑΒ ἡμίσεια ἡ ὑπὸ ΕΑΒ, τῆς  δὲ ὑπὸ ΑΒΓ ἡμίσεια ἡ ὑπὸ
ΕΒΑ, καὶ ἡ ὑπὸ ΕΑΒ ἄρα τῇ ὑπὸ ΕΒΑ ἐστιν ἴση; ὥστε
καὶ πλευρὰ ἡ ΕΑ τῇ ΕΒ ἐστιν ἴση. ὁμοίως δὴ δείχομεν,
ὅτι καὶ ἑκατέρα τῶν ΕΑ, ΕΒ [εὐθειῶν] ἑκατέρᾳ τῶν ΕΓ, ΕΔ
ἴση ἐστίν. αἱ τέσσαρες ἄρα αἱ ΕΑ,  ΕΒ, ΕΓ, ΕΔ ἴσαι ἀλλήλαις
εἰσίν. ὁ ἄρα κέντρῳ τῷ Ε καὶ διαστήματι ἑνὶ τῶν Α, Β, Γ, Δ
κύκλος γραφόμενος ἥχει καὶ διὰ τῶν λοιπῶν σημείων καὶ
ἔσται περιγεγραμμένος περὶ τὸ ΑΒΓΔ τετράγωνον. περιγεγράφθω
ὡς ὁ ΑΒΓΔ.}

{\greekfont Περὶ τὸ δοθὲν ἄρα τετράγωνον κύκλος περιγέγραπται; ὅπερ ἔδει ποιῆσαι.}}

\ParallelRText{
\begin{center}
{\large Proposition 9}
\end{center}

To  circumscribe a circle about a given square.

Let $ABCD$ be the given square. So it is required to circumscribe a circle about
square $ABCD$.

$AC$ and $BD$ being joined, let them cut one another at $E$.\\

\epsfysize=2.2in
\centerline{\epsffile{Book04/fig06e.eps}}

And since $DA$ is equal to $AB$, and $AC$ (is) common, the two (straight-lines)
$DA$, $AC$ are thus equal to the two (straight-lines) $BA$, $AC$. And the base $DC$
(is) equal to the base $BC$. Thus, angle $DAC$ is equal to angle $BAC$ [Prop.~1.8]. Thus, the angle $DAB$ has been cut in half by $AC$. So, similarly,
we can show that $ABC$, $BCD$, and $CDA$ have each been cut in half
by the straight-lines $AC$ and $DB$. And since angle $DAB$ is equal
 to $ABC$, and $EAB$ is half of $DAB$, and $EBA$ half of $ABC$, $EAB$ is thus
 also equal to $EBA$. So that side $EA$ is also equal to  $EB$ [Prop.~1.6]. So, similarly, we can show that each of the [straight-lines] $EA$ and $EB$ are also equal to each of $EC$ and $ED$. Thus, the four (straight-lines)
 $EA$, $EB$, $EC$, and $ED$ are equal to one another. Thus, the circle drawn with
 center $E$, and radius one of $A$, $B$, $C$, or $D$, will also go through the remaining points, and will be circumscribed about the square $ABCD$. Let
 it be (so) circumscribed, like $ABCD$ (in the figure).
 
 Thus, a circle has been circumscribed about the given square.
 (Which is) the very thing it was required to do.}
\end{Parallel}

%%%%%%
% Prop 4.10
%%%%%%
\pdfbookmark[1]{Proposition 4.10}{pdf4.10}
\begin{Parallel}{}{} 
\ParallelLText{
\begin{center}
{\large \ggn{10}.}
\end{center}\vspace*{-7pt}

{\greekfont Ἰσοσκελὲς τρίγωνον συστήσασθαι ἔθον ἑκατέραν τῶν πρὸς τῇ
βάσει γωνιῶν διπλασίονα τῆς λοιπῆς.}

\epsfysize=2.2in
\centerline{\epsffile{Book04/fig10g.eps}}

{\greekfont Ἐκκείσθω τις εὐθεῖα ἡ ΑΒ, καὶ τετμήσθω κατὰ τὸ Γ σημεῖον,
ὥστε τὸ ὑπὸ τῶν ΑΒ, ΒΓ περιεχόμενον ὀρθογώνιον ἴσον
εἶναι τῷ ἀπὸ τῆς ΓΑ τετραγώνῳ; καὶ κέντρῳ τῷ Α καὶ
διαστήματι τῷ ΑΒ κύκλος γεγράφθω ὁ ΒΔΕ, καὶ ἐνηρμόσθω
εἰς τὸν ΒΔΕ κύκλον τῇ ΑΓ εὐθείᾳ μὴ μείζονι οὔσῃ τῆς
τοῦ ΒΔΕ κύκλου διαμέτρου ἴση εὐθεῖα ἡ ΒΔ; καὶ ἐπεζεύθθωσαν
αἱ ΑΔ, ΔΓ, καὶ περιγεγράφθω περὶ τὸ ΑΓΔ τρίγωνον κύκλος ὁ
ΑΓΔ.}

{\greekfont Καὶ ἐπεὶ τὸ ὑπὸ τῶν ΑΒ, ΒΓ ἴσον ἐστὶ τῷ ἀπὸ
τῆς ΑΓ, ἴση δὲ ἡ ΑΓ τῇ ΒΔ, τὸ ἄρα ὑπὸ τῶν ΑΒ, ΒΓ ἴσον
ἐστὶ τῷ ἀπὸ τῆς ΒΔ. καὶ ἐπεὶ κύκλου τοῦ ΑΓΔ εἴληπταί
τι σημεῖον ἐκτὸς τὸ Β, καὶ ἀπὸ τοῦ Β πρὸς τὸν ΑΓΔ κύκλον
προσπεπτώκασι δύο εὐθεῖαι αἱ ΒΑ, ΒΔ, καὶ ἡ μὲν α ὐτῶν
τέμνει, ἡ δὲ προσπίπτει, καί ἐστι τὸ ὑπὸ τῶν ΑΒ, ΒΓ ἴσον τῷ ἀπὸ τῆς ΒΔ, ἡ ΒΔ ἄρα ἐφάπτεται τοῦ ΑΓΔ κύκλου.
ἐπεὶ οὖν ἐφάπτεται μὲν ἡ ΒΔ, ἀπὸ δὲ τῆς κατὰ τὸ Δ
ἐπαφῆς διῆκται ἡ ΔΓ, ἡ ἄρα ὑπὸ ΒΔΓ γωνιά ἴση ἐστὶ
τῇ ἐν τῷ ἐναλλὰχ τοῦ κύκλου τμήματι γωνίᾳ
τῇ ὑπὸ ΔΑΓ. ἐπεὶ οὖν ἴση ἐστὶν ἡ ὑπὸ ΒΔΓ τῇ
ὑπὸ ΔΑΓ, κοινὴ προσκείσθω ἡ ὑπὸ ΓΔΑ; ὅλη ἄρα ἡ
ὑπὸ ΒΔΑ ἴση ἐστὶ δυσὶ ταῖς ὑπὸ ΓΔΑ, ΔΑΓ. ἀλλὰ ταῖς
ὑπὸ ΓΔΑ, ΔΑΓ ἴση ἐστὶν ἡ ἐκτὸς ἡ ὑπὸ ΒΓΔ;
καὶ ἡ ὑπὸ ΒΔΑ ἄρα ἴση ἐστὶ τῇ ὑπὸ ΒΓΔ. ἀλλὰ
ἡ ὑπὸ ΒΔΑ τῇ ὑπὸ ΓΒΔ ἐστιν ἴση, ἐπεὶ καὶ πλευρὰ
ἡ ΑΔ τῇ ΑΒ ἐστιν ἴση; ὥστε καὶ ἡ ὑπὸ ΔΒΑ τῇ
ὑπὸ ΒΓΔ ἐστιν ἴση.
αἱ τρεῖς ἄρα αἱ ὑπὸ ΒΔΑ,
ΔΒΑ, ΒΓΑ ἴσαι ἀλλήλαις εἰσίν. καὶ ἐπεὶ ἴση ἐστὶν
ἡ ὑπὸ ΔΒΓ γωνία τῇ ὑπὸ ΒΓΔ, ἴση ἐστὶ καὶ πλευρὰ
ἡ ΒΔ πλευρᾷ τῇ ΔΓ. ἀλλὰ ἡ ΒΔ τῇ ΓΑ ὑπόκειται ἴση;
καὶ ἡ ΓΑ ἄρα τῇ ΓΔ ἐστιν ἴση; ὥστε καὶ
γωνία ἡ ὑπὸ ΓΔΑ γωνίᾳ τῇ ὑπὸ ΔΑΓ ἐστιν ἴση; αἱ
ἄρα ὑπὸ ΓΔΑ, ΔΑΓ τῆς ὑπὸ ΔΑΓ εἰσι διπλασίους.
ἴση δὲ ἡ ὑπὸ ΒΓΔ ταῖς ὑπὸ ΓΔΑ, ΔΑΓ; καὶ
ἡ ὑπὸ ΒΓΔ ἄρα τῆς ὑπὸ ΓΑΔ ἐστι διπλῆ. ἴση
δὲ ἡ ὑπὸ ΒΓΔ ἑκατέρᾳ τῶν ὑπὸ ΒΔΑ, ΔΒΑ; καὶ
ἑκατέρα ἄρα τῶν ὑπὸ ΒΔΑ, ΔΒΑ τῆς ὑπὸ ΔΑΒ
ἐστι διπλῆ.}

{\greekfont Ἰσοσκελὲς ἄρα τρίγωνον συνέσταται τὸ ΑΒΔ ἔθον
ἑκατέραν τῶν πρὸς τῇ ΔΒ βάσει γωνιῶν διπλασίονα τῆς
λοιπῆς; ὅπερ ἔδει ποιῆσαι.}}

\ParallelRText{
\begin{center}
{\large Proposition 10}
\end{center}

To construct an isosceles triangle having each of the angles at the base
double the remaining (angle).

\epsfysize=2.2in
\centerline{\epsffile{Book04/fig10e.eps}}

Let some straight-line $AB$ be taken, and let it be cut at point $C$
so that the rectangle contained by $AB$ and $BC$ is equal to the square
on $CA$ [Prop.~2.11]. And let the circle $BDE$ be drawn with
center $A$, and radius $AB$. And let the straight-line $BD$, equal to the straight-line $AC$, being not greater than the diameter of circle $BDE$, be inserted into circle $BDE$ [Prop.~4.1]. And let $AD$ and $DC$ be joined. 
And let the circle $ACD$ be circumscribed about triangle $ACD$ [Prop.~4.5].

And since the (rectangle contained) by $AB$ and $BC$ is equal to the
(square) on $AC$, and $AC$ (is) equal to $BD$, the (rectangle contained) by
$AB$ and $BC$ is thus equal to the (square) on $BD$. And since some point
$B$ has been taken outside of circle $ACD$, and two straight-lines $BA$ and
$BD$ have radiated from $B$ towards the circle $ACD$, and (one) of them
cuts (the circle), and (the other) meets (the circle), and the (rectangle contained)
by $AB$ and $BC$ is equal to the (square) on  $BD$,   $BD$ thus touches
circle $ACD$ [Prop.~3.37]. Therefore, since $BD$ touches (the circle),
and $DC$ has been drawn across (the circle) from the point of contact $D$,
the angle $BDC$ is thus equal to the angle $DAC$ in the alternate segment
of the circle [Prop.~3.32]. Therefore, since $BDC$ is equal to
$DAC$, let $CDA$ be added to both. Thus, the whole of $BDA$ is equal
to the two (angles) $CDA$ and $DAC$. But, the
external (angle)  $BCD$  is equal to $CDA$ and $DAC$  [Prop.~1.32]. Thus, $BDA$ is also equal to
$BCD$. But, $BDA$ is equal to $CBD$, since the side $AD$ is also equal to
$AB$ [Prop.~1.5]. So that $DBA$ is also equal to $BCD$. Thus, the three
(angles) $BDA$, $DBA$, and $BCD$ are equal to one another. And since angle
$DBC$ is equal to $BCD$, side $BD$ is also equal to side $DC$ [Prop.~1.6].
But, $BD$ was assumed (to be) equal to $CA$. Thus, $CA$ is also equal to $CD$.
So that angle $CDA$ is also equal to angle $DAC$ [Prop.~1.5]. 
Thus, $CDA$ and $DAC$ is double $DAC$. But $BCD$ (is) equal to $CDA$ and
$DAC$. Thus, $BCD$ is also double $CAD$. And $BCD$ (is) equal to to each of
$BDA$ and $DBA$. Thus, $BDA$ and $DBA$ are each double $DAB$.

Thus, the isosceles triangle $ABD$ has been constructed having each of the
angles at the base $BD$ double the remaining (angle). (Which is) the
very thing it was required to do.}
\end{Parallel}

%%%%%%
% Prop 4.11
%%%%%%
\pdfbookmark[1]{Proposition 4.11}{pdf4.11}
\begin{Parallel}{}{} 
\ParallelLText{
\begin{center}
{\large \ggn{11}.}
\end{center}\vspace*{-7pt}

{\greekfont Εἰς τὸν δοθέντα κύκλον πεντάγωνον ἰσόπλευρόν τε καὶ ἰσογώνιον
ἐγγράψαι.}

\epsfysize=1.8in
\centerline{\epsffile{Book04/fig11g.eps}}

{\greekfont Ἔστω ὁ δοθεὶς κύκλος ὁ ΑΒΓΔΕ; δεῖ δὴ εἰς τὸν ΑΒΓΔΕ κύκλον πεντάγωνον ἰσόπλευρόν τε καὶ ἰσογώνιον ἐγγράψαι.}

{\greekfont Ἐκκείσθω τρίγωνον ἰσοσκελὲς τὸ ΖΗΘ διπλασίονα ἔθον ἑκατέραν
τῶν πρὸς τοῖς Η, Θ γωνιῶν τῆς πρὸς τῷ Ζ, καὶ ἐγγεγράφθω
εἰς τὸν ΑΒΓΔΕ κύκλον τῷ ΖΗΘ τριγώνῳ ἰσογώνιον τρίγωνον
τὸ ΑΓΔ, ὥστε τῇ μὲν πρὸς τῷ Ζ γωνίᾳ ἴσην εἶναι
τὴν ὑπὸ ΓΑΔ, ἑκατέραν δὲ τῶν πρὸς τοῖς Η, Θ ἴσην ἑκατέρᾳ
τῶν ὑπὸ ΑΓΔ, ΓΔΑ; καὶ ἑκατέρα ἄρα τῶν ὑπὸ ΑΓΔ, ΓΔΑ
τῆς ὑπὸ ΓΑΔ ἐστι διπλῆ. τετμήσθω δὴ ἑκατέρα τῶν ὑπὸ
ΑΓΔ, ΓΔΑ δίθα ὑπὸ ἑκατέρας τῶν ΓΕ, ΔΒ εὐθειῶν, καὶ
ἐπεζεύθθωσαν αἱ ΑΒ, ΒΓ,  ΔΕ, ΕΑ.}

{\greekfont Ἐπεὶ οὖν ἑκατέρα τῶν ὑπὸ ΑΓΔ, ΓΔΑ γωνιῶν διπλασίων
ἐστὶ τῆς ὑπὸ ΓΑΔ, καὶ τετμ ημέναι εἰσὶ δίθα ὑπὸ τῶν ΓΕ,
ΔΒ εὐθειῶν, αἱ πέντε ἄρα γωνίαι αἱ ὑπὸ ΔΑΓ, ΑΓΕ, ΕΓΔ, ΓΔΒ,
ΒΔΑ ἴσαι ἀλλήλαις εἰσίν. αἱ δὲ ἴσαι γωνίαι ἐπὶ ἴσων
περιφερειῶν βεβήκασιν; αἱ πέντε  ἄρα περιφέρειαι αἱ ΑΒ, ΒΓ, ΓΔ, ΔΕ, ΕΑ ἴσαι ἀλλήλαις εἰσίν. ὑπὸ δὲ τὰς ἴσας
περιφερείας ἴσαι εὐθεῖαι ὑποτείνουσιν; αἱ πέντε ἄρα εὐθεῖαι
αἱ ΑΒ, ΒΓ, ΓΔ, ΔΕ, ΕΑ ἴσαι ἀλλήλαις εἰσίν; ἰσόπλευρον ἄρα
ἐστὶ τὸ ΑΒΓΔΕ πεντάγωνον. λέγω δή, ὅτι καὶ ἰσογώνιον. ἐπεὶ
γὰρ ἡ ΑΒ περιφέρεια τῇ ΔΕ περιφερείᾳ ἐστὶν ἴση, κοινὴ
προσκείσθω ἡ ΒΓΔ; ὅλη ἄρα ἡ ΑΒΓΔ περιφέρια  ὅλῃ
τῇ ΕΔΓΒ
περιφερείᾳ ἐστὶν ἴση. καὶ βέβηκεν ἐπὶ μὲν τῆς ΑΒΓΔ
περιφερείας γωνία ἡ ὑπὸ ΑΕΔ, ἐπὶ δὲ τῆς ΕΔΓΒ περιφερείας
γωνία ἡ ὑπὸ ΒΑΕ; καὶ ἡ ὑπὸ ΒΑΕ ἄρα γωνία
τῇ ὑπὸ ΑΕΔ ἐστιν ἴση. διὰ τὰ α ὐτὰ δὴ καὶ ἑκάστη τῶν
ὑπὸ ΑΒΓ, ΒΓΔ, ΓΔΕ γωνιῶν ἑκατέρᾳ τῶν ὑπὸ ΒΑΕ, ΑΕΔ
ἐστιν ἴση; ἰσογώνιον ἄρα ἐστὶ τὸ ΑΒΓΔΕ πεντάγωνον.
ἐδείθθη δὲ καὶ ἰσόπλευρον.}

{\greekfont Εἰς ἄρα τὸν δοθέντα κύκλον πεντάγωνον ἰσόπλευρόν τε καὶ ἰσογώνιον
ἐγγέγραπται; ὅπερ ἔδει ποιῆσαι.}}

\ParallelRText{
\begin{center}
{\large Proposition 11}
\end{center}

To inscribe an equilateral and equiangular pentagon in a given circle.

\epsfysize=1.8in
\centerline{\epsffile{Book04/fig11e.eps}}

Let $ABCDE$ be the given circle. So it is required to inscribed an equilateral
and equiangular pentagon in circle $ABCDE$.

Let the the isosceles triangle $FGH$ be set up having each of the angles
at $G$ and $H$ double the (angle) at $F$ [Prop.~4.10]. And let triangle $ACD$, equiangular
to $FGH$, be inscribed in circle $ABCDE$, such that $CAD$ is equal 
to the angle at $F$, and  the (angles) at $G$ and $H$ (are) equal to $ACD$ and $CDA$, respectively [Prop.~4.2]. Thus,
$ACD$ and $CDA$ are each double $CAD$. So let $ACD$ and $CDA$ 
be cut in half by the straight-lines $CE$ and $DB$, respectively
[Prop.~1.9]. And let $AB$, $BC$, $DE$ and
$EA$ be joined.

Therefore, since angles $ACD$ and $CDA$ are each double $CAD$, and are cut
in half by the straight-lines $CE$ and $DB$, the five angles
$DAC$, $ACE$, $ECD$, $CDB$, and $BDA$ are thus equal to one another.
And equal angles stand upon equal circumferences [Prop.~3.26]. Thus, the five circumferences
$AB$, $BC$, $CD$, $DE$, and $EA$ are equal to one another [Prop.~3.29]. Thus, the pentagon
$ABCDE$ is equilateral. So I say that (it is) also equiangular. For since the
circumference $AB$ is equal to the circumference $DE$, let $BCD$ be added to both. Thus, the whole circumference $ABCD$ is equal to the
whole circumference $EDCB$. And the angle $AED$ stands upon circumference
$ABCD$, and angle $BAE$ upon circumference $EDCB$. Thus, angle $BAE$ is
also equal to $AED$ [Prop.~3.27]. 
So, for the same (reasons), each of the angles $ABC$, $BCD$,  and $CDE$ is also equal
to each of $BAE$ and $AED$. Thus, pentagon $ABCDE$ is equiangular.
And it was also shown (to be) equilateral.

Thus, an equilateral and equiangular pentagon has been inscribed in the
given circle. (Which is) the very thing it was required to do.}
\end{Parallel}

%%%%%%
% Prop 4.12
%%%%%%
\pdfbookmark[1]{Proposition 4.12}{pdf4.12}
\begin{Parallel}{}{} 
\ParallelLText{
\begin{center}
{\large \ggn{12}.}
\end{center}\vspace*{-7pt}

{\greekfont Περὶ τὸν δοθέντα κύκλον πεντάγωνον ἰσόπλευρόν τε καὶ ἰσογώνιον
περιγράψαι.}

\epsfysize=2in
\centerline{\epsffile{Book04/fig12g.eps}}

{\greekfont Ἔστω ὁ δοθεὶς κύκλος ὁ ΑΒΓΔΕ; δεῖ δὲ περὶ τὸν ΑΒΓΔΕ κύκλον
πεντάγωνον ἰσόπλευρόν τε καὶ ἰσογώνιον περιγράψαι.}

{\greekfont Νενοήσθω τοῦ ἐγγεγραμμένου πενταγώνου τῶν γωνιῶν σημεῖα
τὰ Α, Β, Γ, Δ, Ε, ὥστε ἴσας εἶναι τὰς ΑΒ, ΒΓ, ΓΔ, ΔΕ, ΕΑ
περιφερείας; καὶ διὰ τῶν Α, Β, Γ, Δ, Ε ἤθθωσαν τοῦ κύκλου
ἐφαπτόμεναι αἱ ΗΘ, ΘΚ, ΚΛ, ΛΜ, μ η, καὶ εἰλήφθω τοῦ ΑΒΓΔΕ
κύκλου κέντρον τὸ Ζ, καὶ ἐπεζεύθθωσαν αἱ ΖΒ, ΖΚ, ΖΓ, ΖΛ, ΖΔ.}

{\greekfont Καὶ ἐπεὶ ἡ μὲν ΚΛ εὐθεῖα ἐφάπτεται τοῦ ΑΒΓΔΕ κατὰ τὸ
Γ, ἀπὸ δὲ τοῦ Ζ κέντρου ἐπὶ τὴν κατὰ τὸ Γ ἐπαφὴν ἐπέζευκται
ἡ ΖΓ, ἡ ΖΓ ἄρα κάθετός ἐστιν ἐπὶ τὴν ΚΛ; ὀρθὴ ἄρα ἐστὶν
ἑκατέρα τῶν πρὸς τῷ Γ γωνιῶν. διὰ τὰ α ὐτὰ δὴ καὶ αἱ 
πρὸς τοῖς Β, Δ σημείοις γωνίαι  ὀρθαί εἰσιν. καὶ ἐπεὶ ὀρθή 
ἐστιν ἡ ὑπὸ ΖΓΚ γωνία, τὸ ἄρα ἀπὸ τῆς ΖΚ ἴσον 
ἐστὶ τοῖς ἀπὸ τῶν ΖΓ, ΓΚ. διὰ τὰ α ὐτὰ δὴ καὶ τοῖς 
ἀπὸ τῶν ΖΒ, ΒΚ ἴσον ἐστὶ τὸ ἀπὸ τῆς ΖΚ; ὥστε τὰ ἀπὸ
τῶν ΖΓ, ΓΚ τοῖς ἀπὸ τῶν ΖΒ, ΒΚ ἐστιν ἴσα, ὧν τὸ
ἀπὸ τῆς ΖΓ τῷ ἀπὸ τῆς ΖΒ ἐστιν ἴσον; λοιπὸν ἄρα
τὸ ἀπὸ τῆς ΓΚ τῷ ἀπὸ τῆς ΒΚ ἐστιν ἴσον. ἴση
ἄρα ἡ ΒΚ τῇ ΓΚ. καὶ ἐπεὶ ἴση ἐστὶν ἡ ΖΒ τῇ ΖΓ, καὶ κοινὴ 
ἡ ΖΚ, δύο δὴ αἱ ΒΖ, ΖΚ δυσὶ ταῖς ΓΖ, ΖΚ ἴσαι εἰσίν;
καὶ βάσις ἡ ΒΚ βάσει τῇ ΓΚ [ἐστιν] ἴση; γωνία ἄρα ἡ μὲν 
ὑπὸ ΒΖΚ [γωνίᾳ] τῇ ὑπὸ ΚΖΓ ἐστιν ἴση; ἡ δὲ ὑπὸ ΒΚΖ
τῇ ὑπὸ ΖΚΓ; διπλῆ ἄρα ἡ μὲν ὑπὸ ΒΖΓ τῆς ὑπὸ ΚΖΓ,
ἡ δὲ ὑπὸ ΒΚΓ τῆς ὑπὸ ΖΚΓ. διὰ τὰ α ὐτὰ δὴ καὶ ἡ
μὲν ὑπὸ ΓΖΔ τῆς ὑπὸ ΓΖΛ ἐστι διπλῆ, ἡ δὲ ὑπὸ ΔΛΓ
τῆς ὑπὸ ΖΛΓ. καὶ ἐπεὶ ἴση ἐστὶν ἡ ΒΓ περιφέρεια τῇ ΓΔ,
ἴση ἐστὶ καὶ γωνία ἡ ὑπὸ ΒΖΓ τῇ ὑπὸ ΓΖΔ.
καί ἐστιν ἡ μὲν ὑπὸ ΒΖΓ τῆς ὑπὸ ΚΖΓ διπλῆ,
ἡ δὲ ὑπὸ ΔΖΓ τῆς ὑπὸ ΛΖΓ; ἴση ἄρα καὶ ἡ ὑπὸ
ΚΖΓ τῇ ὑπὸ ΛΖΓ; ἐστὶ δὲ καὶ ἡ ὑπὸ ΖΓΚ γωνία
τῇ ὑπὸ ΖΓΛ ἴση. δύο δὴ τρίγωνά ἐστι τὰ ΖΚΓ,
ΖΛΓ τὰς  δύο γωνίας ταῖς δυσὶ γωνίαις ἴσας ἔθοντα καὶ
μίαν πλευρὰν μιᾷ πλευρᾷ ἴσην κοινὴν α ὐτῶν
τὴν ΖΓ; καὶ τὰς λοιπὰς ἄρα πλευρὰς ταῖς λοιπαῖς πλευραῖς ἴσας
ἕχει καὶ τὴν λοιπὴν γωνίαν τῇ λοιπῇ γωνίᾳ; ἴση ἄρα ἡ μὲν ΚΓ εὐθεῖα τῇ ΓΛ, ἡ δὲ ὑπὸ
ΖΚΓ γωνία τῇ ὑπὸ ΖΛΓ. καὶ ἐπεὶ ἴση ἐστὶν
ἡ ΚΓ τῇ ΓΛ, διπλῆ ἄρα ἡ ΚΛ τῆς ΚΓ. διὰ τὰ α ὐτα
δὴ δειθθήσεται καὶ ἡ 
ΘΚ τῆς ΒΚ διπλῆ. καί ἐστιν ἡ ΒΚ τῇ ΚΓ
ἴση; καὶ ἡ ΘΚ ἄρα τῇ ΚΛ ἐστιν ἴση. ὁμοίως δὴ δειθθήσεται
καὶ ἑκάστη τῶν ΘΗ, ΗΜ,
ΜΛ ἑκατέρᾳ τῶν ΘΚ, ΚΛ ἴση; ἰσόπλευρον ἄρα ἐστὶ τὸ ΗΘΚΛΜ
πεντάγωνον. λέγω δή, ὅτι καὶ ἰσογώνιον. ἐπεὶ γὰρ ἴση
ἐστὶν ἡ ὑπὸ ΖΚΓ γωνία τῇ ὑπὸ ΖΛΓ, καὶ ἐδείθθη τῆς
μὲν ὑπὸ ΖΚΓ διπλῆ ἡ ὑπὸ
ΘΚΛ, τῆς δὲ ὑπὸ ΖΛΓ διπλῆ ἡ ὑπὸ ΚΛΜ, καὶ ἡ ὑπὸ
ΘΚΛ ἄρα τῇ ὑπὸ ΚΛΜ ἐστιν ἴση. ὁμοίως δὴ δειθθήσεται καὶ
ἑκάστη τῶν ὑπὸ ΚΘΗ, ΘΗΜ, ΗΜΛ ἑκατέρᾳ τῶν ὑπὸ ΘΚΛ, ΚΛΜ
ἴση; αἱ πέντε ἄρα γωνίαι αἱ ὑπὸ ΗΘΚ, ΘΚΛ, ΚΛΜ, Λμ η,
μ ηΘ
ἴσαι ἀλλήλαις εἰσίν. ἰσογώνιον ἄρα ἐστὶ τὸ ΗΘΚΛΜ
πεντάγωνον. ἐδείθθη δὲ καὶ ἰσόπλευρον, καὶ
περιγέγραπται περὶ τὸν ΑΒΓΔΕ κύκλον.}

{\greekfont μ̀βοχ{[}Περὶ τὸν δοθέντα ἄρα κύκλον πεντάγωνον ἰσόπλευρ-όν
τε καὶ ἰσογώνιον περιγέγραπται]; ὅπερ ἔδει
ποιῆσαι.}}

\ParallelRText{
\begin{center}
{\large Proposition 12}
\end{center}

To circumscribe an equilateral and equiangular
pentagon about a given circle.

\epsfysize=2in
\centerline{\epsffile{Book04/fig12e.eps}}

Let $ABCDE$ be the given circle. So it is required to circumscribe an
equilateral and equiangular pentagon about circle $ABCDE$. 

Let $A$, $B$, $C$, $D$,  and $E$ be conceived as the angular points of a
pentagon having been inscribed (in circle $ABCDE$) [Prop.~3.11], such that the circumferences $AB$, $BC$, $CD$, $DE$, and
$EA$ are equal. And let $GH$, $HK$, $KL$, $LM$, and $MG$  be drawn through
(points) $A$, $B$, $C$, $D$, and $E$ (respectively), touching the circle.$^\dag$ And let the center $F$
of the circle $ABCDE$ be found [Prop.~3.1].
And let $FB$, $FK$, $FC$, $FL$, and $FD$ be joined.

And since the straight-line $KL$ touches (circle) $ABCDE$ at $C$, and $FC$ has been
joined from the center $F$ to the point of contact $C$, $FC$ is thus perpendicular 
to $KL$ [Prop.~3.18]. Thus, each of the angles
at $C$ is a right-angle. So, for the same (reasons), the angles at $B$ and $D$
are also right-angles. And since angle $FCK$ is a right-angle, the (square)
on $FK$ is thus equal to the (sum of the squares) on $FC$ and $CK$ [Prop.~1.47]. So, for the same (reasons), 
the (square) on $FK$ is also equal to the (sum of the squares) on $FB$ and $BK$.
So that the (sum of the squares) on $FC$ and $CK$ is equal to the (sum of the squares) on $FB$ and $BK$, of which the (square) on $FC$ is equal to the
(square) on $FB$. Thus, the remaining (square) on $CK$ is equal to
the remaining (square) on $BK$. Thus, $BK$ (is) equal to $CK$. And since
$FB$ is equal to $FC$, and $FK$ (is) common, the two (straight-lines)
$BF$, $FK$ are equal to the two (straight-lines) $CF$, $FK$. And the base $BK$
[is] equal to the base $CK$. Thus, angle $BFK$ is equal to [angle] $KFC$ 
[Prop.~1.8]. And $BKF$ (is equal) to $FKC$ [Prop.~1.8].
Thus, $BFC$ (is) double $KFC$,
and $BKC$ (is double) $FKC$. So, for the same (reasons), $CFD$ is also double
$CFL$, and $DLC$ (is also double) $FLC$. And since circumference $BC$
is equal  to $CD$, angle $BFC$ is also  equal  to $CFD$ [Prop.~3.27]. And $BFC$ is double $KFC$,
and $DFC$ (is double) $LFC$.  Thus, $KFC$ is also equal to $LFC$. 
And
angle $FCK$ is also equal to $FCL$. So, $FKC$ and $FLC$ are two triangles
having two angles equal to two angles, and one side equal to one side,
(namely) their common (side) $FC$. Thus, they will also have the
remaining sides equal to the (corresponding) remaining sides, and
the remaining angle to the remaining angle [Prop.~1.26]. Thus, the straight-line $KC$
(is) equal to $CL$, and the angle $FKC$ to $FLC$. And since $KC$ is
equal to $CL$, $KL$ (is) thus double $KC$. So, for the
same (reasons), it can be shown that $HK$ (is) also double $BK$. And
$BK$ is equal to $KC$. Thus, $HK$ is also equal to $KL$. So, similarly,  each of $HG$, $GM$, and $ML$ can also be shown (to be) equal to each of
$HK$ and $KL$. Thus, pentagon $GHKLM$ is equilateral. So I say that
(it is) also equiangular. For since angle $FKC$ is equal to $FLC$, 
and $HKL$ was shown (to be) double $FKC$, and $KLM$ double 
$FLC$, $HKL$ is thus also equal to $KLM$. So, similarly, each of $KHG$, $HGM$, and $GML$ can also be shown (to be) equal to each of
$HKL$ and $KLM$. Thus, the five angles $GHK$, $HKL$, $KLM$, $LMG$, and
$MGH$ are equal to one another. Thus, the pentagon $GHKLM$ is
equiangular. And it was also shown (to be) equilateral, and has been circumscribed 
about circle $ABCDE$.

\mbox{[}Thus, an equilateral and equiangular pentagon has been circumscribed
about the given circle]. (Which is) the very thing it was required to do. }
\end{Parallel}
{\footnotesize \noindent$^\dag$ See the footnote to Prop.~3.34.} 

%%%%%%
% Prop 4.13
%%%%%%
\pdfbookmark[1]{Proposition 4.13}{pdf4.13}
\begin{Parallel}{}{} 
\ParallelLText{
\begin{center}
{\large \ggn{13}.}
\end{center}\vspace*{-7pt}

{\greekfont Εἰς τὸ δοθὲν πεντάγωνον, ὅ ἐστιν ἰσόπλευρόν τε καὶ ἰσογώνιον, κύκλον ἐγγράψαι.}

{\greekfont Ἔστω τὸ δοθὲν πεντάγωνον ἰσόπλευρόν τε καὶ ἰσογώνι-ον
τὸ ΑΒΓΔΕ; δεῖ δὴ εἰς τὸ ΑΒΓΔΕ πεντάγωνον κύκλον ἐγγράψαι.}

{\greekfont Τετμήσθω γὰρ ἑκατέρα τῶν ὑπὸ ΒΓΔ, ΓΔΕ γωνιῶν δίθα ὑπὸ
ἑκατέρας τῶν ΓΖ, ΔΖ εὐθειῶν; καὶ ἀπὸ τοῦ Ζ σημείου, καθ ὃ
συμβάλλουσιν ἀλλήλαις αἱ ΓΖ, ΔΖ εὐθεῖαι, ἐπεζεύθθωσαν αἱ ΖΒ,
ΖΑ, ΖΕ εὐθεῖαι. καὶ ἐπεὶ ἴση ἐστὶν ἡ ΒΓ τῇ ΓΔ, κοινὴ
δὲ ἡ ΓΖ, δύο δὴ αἱ ΒΓ, ΓΖ δυσὶ ταῖς ΔΓ, ΓΖ ἴσαι
εἰσίν; καὶ γωνία ἡ ὑπὸ ΒΓΖ γωνίᾳ τῇ ὑπὸ ΔΓΖ [ἐστιν]
ἴση; βάσις ἄρα ἡ ΒΖ βάσει τῇ ΔΖ ἐστιν ἴση, 
καὶ τὸ ΒΓΖ τρίγωνον τῷ ΔΓΖ τριγώνῳ ἐστιν ἴσον,
καὶ αἱ λοιπαὶ
γωνίαι ταῖς λοιπαῖς γωνίαις ἴσαι ἔσονται, ὑφ ἃς αἱ
ἴσαι πλευραὶ ὑποτείνουσιν; ἴση ἄρα ἡ ὑπὸ ΓΒΖ γωνία 
τῇ  ὑπὸ ΓΔΖ. καὶ ἐπεὶ διπλῆ ἐστιν ἡ ὑπὸ ΓΔΕ τῆς 
ὑπὸ ΓΔΖ, ἴση δὲ ἡ μὲν ὑπὸ ΓΔΕ τῇ ὑπὸ ΑΒΓ, ἡ δὲ 
ὑπὸ ΓΔΖ τῇ ὑπὸ ΓΒΖ, καὶ ἡ ὑπὸ ΓΒΑ ἄρα τῆς ὑπὸ
ΓΒΖ ἐστι διπλῆ; ἴση ἄρα ἡ ὑπὸ ΑΒΖ γωνία τῇ ὑπὸ
ΖΒΓ; ἡ ἄρα ὑπὸ ΑΒΓ γωνία δίθα τέτμ ηται ὑπὸ τῆς
ΒΖ εὐθείας. ὁμοίως δὴ δειθθήσεται, ὅτι καὶ ἑκατέρα τῶν
ὑπὸ ΒΑΕ, ΑΕΔ δίθα τέτμ ηται ὑπὸ ἑκατέρας τῶν ΖΑ, ΖΕ
εὐθειῶν. ἤθθωσαν δὴ ἀπὸ τοῦ Ζ σημείου ἐπὶ τὰς ΑΒ, ΒΓ,
ΓΔ, ΔΕ, ΕΑ εὐθείας κάθετοι αἱ ΖΗ, ΖΘ, ΖΚ, ΖΛ, ΖΜ. καὶ ἐπεὶ
ἴση ἐστὶν ἡ ὑπὸ ΘΓΖ γωνία τῇ ὑπὸ ΚΓΖ, ἐστὶ δὲ 
καὶ ὀρθὴ ἡ ὑπὸ ΖΘΓ
[ὀρθῇ] τῇ ὑπὸ ΖΚΓ ἴση,
δύο δὴ τρίγωνά ἐστι τὰ ΖΘΓ, ΖΚΓ τὰς δύο γωνίας δυσὶ γωνίαις ἴσας
ἔθοντα καὶ μίαν πλευρὰν μιᾷ πλευρᾷ ἴσην κοινὴν α ὐτῶν 
τὴν ΖΓ ὑποτείνουσαν ὑπὸ μίαν τῶν ἴσων γωνιῶν; καὶ
τὰς λοιπὰς ἄρα πλευρὰς ταῖς λοιπαῖς πλευραῖς ἴσας ἕχει;
ἴση ἄρα ἡ ΖΘ κάθετος τῂ ΖΚ καθέτῳ. ὁμοίως
δὴ δειθθήσεται, ὅτι καὶ ἑκάστη τῶν ΖΛ, ΖΜ, ΖΗ
ἑκατέρᾳ τῶν ΖΘ, ΖΚ ἴση ἐστίν; αἱ πέντε ἄρα εὐθεῖαι αἱ
ΖΗ, ΖΘ, ΖΚ, ΖΛ, ΖΜ  ἴσαι ἀλλήλαις εἰσίν. ὁ ἄρα
κέντρῳ τῷ Ζ διαστήματι δὲ ἑνὶ τῶν Η, Θ, Κ, Λ, Μ
κύκλος γραφόμενος ἥχει καὶ διὰ τῶν λοιπῶν σημείων
καὶ ἐφάψεται τῶν 
ΑΒ, ΒΓ, ΓΔ, ΔΕ, ΕΑ εὐθειῶν διὰ τὸ ὀρθὰς εἶναι
τὰς πρὸς τοῖς Η, Θ, Κ, Λ, Μ σημείοις γωνίας. εἰ γὰρ οὐκ
ἐφάψεται α ὐτῶν, ἀλλὰ τεμεῖ α ὐτάς, συμβήσεται τὴν τῇ
διαμέτρῳ τοῦ κύκλου πρὸς ὀρθὰς ἀπ ἄκρας ἀγομένην
ἐντὸς πίπτειν τοῦ κύκλου; ὅπερ ἄτοπον ἐδείθθη.
οὐκ ἄρα ὁ κέντρῳ τῷ Ζ διαστήματι δὲ ἑνὶ τῶν
Η, Θ, Κ, Λ, Μ σημείων γραφόμενος κύκλος τεμεῖ τὰς ΑΒ, ΒΓ, ΓΔ, ΔΕ,
ΕΑ εὐθείας;  ἐφάψεται ἄρα α ὐτῶν. γεγράφθω ὡς ὁ ΗΘΚΛΜ.}\\~\\~\\

\epsfysize=2.2in
\centerline{\epsffile{Book04/fig13g.eps}}

{\greekfont Εἰς ἄρα τὸ δοθὲν πεντάγωνον, ὅ ἐστιν ἰσόπλευρόν
τε καὶ ἰσογώνιον, κύκλος ἐγγέγραπται; ὅπερ ἔδει ποιῆσαι.}}

\ParallelRText{
\begin{center}
{\large Proposition 13}
\end{center}

To inscribe a circle in a given
pentagon, which is equilateral and equiangular.

Let $ABCDE$ be the given equilateral and equiangular pentagon. So
it is required to inscribe a circle in pentagon $ABCDE$.

For let angles $BCD$ and $CDE$  each be cut in half by each
of the straight-lines $CF$ and $DF$ (respectively) [Prop.~1.9]. And from the point $F$, at which the
straight-lines $CF$ and $DF$ meet one another, let the straight-lines
$FB$, $FA$, and $FE$ be joined. And since $BC$ is equal to $CD$,
and $CF$ (is) common, the two (straight-lines) $BC$, $CF$ 
are equal to the two (straight-lines) $DC$, $CF$. And angle $BCF$ [is]
equal to angle $DCF$. Thus, the base $BF$ is equal to the base $DF$,
and triangle $BCF$ is equal to triangle $DCF$, and the remaining angles
will be equal to the (corresponding) remaining angles which the
equal sides subtend
[Prop.~1.4]. Thus, angle $CBF$ (is) equal
to $CDF$. And since $CDE$ is double $CDF$, and $CDE$ (is) equal to $ABC$,
and $CDF$ to $CBF$, $CBA$ is thus also double $CBF$. Thus, angle $ABF$
is equal to $FBC$. Thus, angle $ABC$ has been cut in half by the straight-line
$BF$. So, similarly, it can be shown that $BAE$ and $AED$ have 
been cut in half by the straight-lines $FA$ and $FE$, respectively.
So let $FG$, $FH$, $FK$, $FL$, and $FM$ be drawn from point $F$,
perpendicular to the straight-lines $AB$, $BC$, $CD$, $DE$, and $EA$ 
(respectively)
[Prop.~1.12]. And since angle
$HCF$ is equal to $KCF$, and the right-angle $FHC$ is also equal to
the [right-angle] $FKC$, $FHC$ and $FKC$ are two triangles having two
angles equal to two angles, and one side equal to one side, (namely)
their common (side) $FC$, subtending one of the equal angles. 
Thus, they will also have the remaining sides equal to the (corresponding)
remaining sides [Prop.~1.26].
Thus, the perpendicular $FH$ (is) equal to the perpendicular $FK$. So,
similarly, it can be shown that $FL$, $FM$, and $FG$ are each equal to
each  of $FH$ and $FK$. Thus, the five straight-lines $FG$, $FH$, $FK$, $FL$,
and $FM$ are equal to one another.
Thus, the circle drawn with center $F$, and
radius one of $G$, $H$, $K$, $L$, or $M$, will also go through the remaining
points, and will touch the straight-lines $AB$, $BC$, $CD$, $DE$, and $EA$,
on account of the angles at points $G$, $H$, $K$, $L$, and $M$ being
right-angles. For if it does not touch them, but cuts them, 
 it follows that a (straight-line) drawn at right-angles to the diameter
 of the circle, from its extremity, falls inside the circle.
 The very thing was shown (to be) absurd [Prop.~3.16]. Thus, the circle drawn
 with center $F$, and radius one of $G$, $H$, $K$, $L$,  or $M$, does not
 cut the straight-lines $AB$, $BC$, $CD$, $DE$, or $EA$. Thus, it will
 touch them. Let it be drawn, like $GHKLM$ (in the figure).

\epsfysize=2.2in
\centerline{\epsffile{Book04/fig13e.eps}}
 
 Thus, a circle has been inscribed in the given pentagon which is
 equilateral and equiangular. (Which is) the very thing it was required to
 do.}
\end{Parallel}

%%%%%%
% Prop 4.14
%%%%%%
\pdfbookmark[1]{Proposition 4.14}{pdf4.14}
\begin{Parallel}{}{} 
\ParallelLText{
\begin{center}
{\large \ggn{14}.}
\end{center}\vspace*{-7pt}

{\greekfont Περὶ τὸ δοθὲν πεντάγωνον, ὅ ἐστιν ἰσόπλευρόν τε καὶ 
ἰσογώνιον, κύκλον περιγράψαι.}

{\greekfont Ἔστω τὸ δοθὲν πεντάγωνον, ὅ ἐστιν ἰσόπλευρόν τε καὶ
ἰσογώνιον, τὸ ΑΒΓΔΕ; δεῖ δὴ περὶ τὸ ΑΒΓΔΕ πεντάγωνον
κύκλον περιγράψαι.}

\epsfysize=2.2in
\centerline{\epsffile{Book04/fig14g.eps}}

{\greekfont Τετμήσθω δὴ ἑκατέρα τῶν ὑπὸ ΒΓΔ, ΓΔΕ γωνιῶν δίθα
ὑπὸ ἑκατέρας τῶν ΓΖ, ΔΖ, καὶ ἀπὸ τοῦ Ζ σημείου,
καθ ὃ συμβάλλουσιν αἱ εὐθεῖαι, ἐπὶ τὰ Β, Α, Ε σημεῖα
ἐπεζεύθθωσαν εὐθεῖαι αἱ ΖΒ, ΖΑ, ΖΕ. ὁμοίως δὴ τῷ πρὸ τούτου
δειθθήσεται, ὅτι καὶ ἑκάστη τῶν ὑπὸ ΓΒΑ, ΒΑΕ, ΑΕΔ
γωνιῶν δίθα τέτμ ηται ὑπὸ ἑκάστης τῶν ΖΒ, ΖΑ, ΖΕ
εὐθειῶν. καὶ ἐπεὶ ἴση ἐστὶν ἡ ὑπὸ ΒΓΔ γωνία τῇ
ὑπὸ ΓΔΕ, καί ἐστι τῆς μὲν ὑπὸ ΒΓΔ  ἡμίσεια ἡ ὑπὸ ΖΓΔ,
τῆς δὲ ὑπὸ ΓΔΕ ἡμίσεια ἡ ὑπὸ ΓΔΖ, καὶ ἡ ὑπὸ
ΖΓΔ ἄρα τῇ ὑπὸ ΖΔΓ ἐστιν ἴση; ὥστε καὶ πλευρὰ ἡ ΖΓ
πλευρᾷ τῇ ΖΔ ἐστιν ἴση. ὁμοίως δὴ δειθθήσεται, ὅτι καὶ
ἑκάστη τῶν ΖΒ, ΖΑ, ΖΕ ἑκατέρᾳ τῶν ΖΓ, ΖΔ ἐστιν ἴση; αἱ
πέντε ἄρα εὐθεῖαι αἱ ΖΑ, ΖΒ, ΖΓ, ΖΔ, ΖΕ  ἴσαι ἀλλήλαις
εἰσίν. ὀ ἄρα κέντρῳ τῷ Ζ καὶ διαστήματι ἑνὶ τῶν ΖΑ, ΖΒ, 
ΖΓ, ΖΔ, ΖΕ κύκλος γραφόμενος ἥχει καὶ διὰ τῶν λοιπῶν
σημείων καὶ ἔσται περιγεγραμμένος. περιγεγράφθω
καὶ ἔστω ὁ ΑΒΓΔΕ.}

{\greekfont Περὶ ἄρα τὸ δοθὲν πεντάγωνον, ὅ ἐστιν ἰσόπλευρόν τε καὶ
ἰσογώνιον, κύκλος περιγέγραπται; ὅπερ ἔδει ποιῆσαι.}}

\ParallelRText{
\begin{center}
{\large Proposition 14}
\end{center}

To circumscribe a circle about a given pentagon which is equilateral and equiangular.

Let $ABCDE$ be the given pentagon  which is equilateral and equiangular.
So it is required to circumscribe a circle about the pentagon 
$ABCDE$.

\epsfysize=2.2in
\centerline{\epsffile{Book04/fig14e.eps}}

So let angles $BCD$ and $CDE$ have  been cut
in half by  the (straight-lines) $CF$ and $DF$, respectively [Prop.~1.9]. And let the straight-lines
$FB$, $FA$, and $FE$ be joined from point $F$, at which the
straight-lines meet, to the points $B$, $A$, and $E$ (respectively).
So, similarly, to the (proposition) before this (one), it can be shown that
angles $CBA$, $BAE$, and $AED$ have also  been cut in half by
the straight-lines $FB$, $FA$, and $FE$, respectively. And since
angle $BCD$ is equal to $CDE$, and $FCD$ is half of $BCD$, and $CDF$ half
of $CDE$, $FCD$ is thus also equal to $FDC$. So that side $FC$ is also equal to
side $FD$ [Prop.~1.6]. So, similarly, it can be
shown that $FB$, $FA$, and $FE$ are also each equal to each of $FC$ and $FD$. 
Thus, the five straight-lines $FA$, $FB$, $FC$, $FD$, and $FE$ are equal to
one another. Thus, the circle drawn with center $F$, and radius one of
$FA$, $FB$, $FC$, $FD$, or $FE$, will also go through the remaining points, and
will be circumscribed. Let it be (so) circumscribed, and let
it be $ABCDE$.

Thus, a circle has been circumscribed about the given pentagon, which
is equilateral and equiangular. (Which is) the very thing it was required to
do.}
\end{Parallel}

%%%%%%
% Prop 4.15
%%%%%%
\pdfbookmark[1]{Proposition 4.15}{pdf4.15}
\begin{Parallel}{}{} 
\ParallelLText{
\begin{center}
{\large \ggn{15}.}
\end{center}\vspace*{-7pt}

{\greekfont Εἰς τὸν δοθέντα κύκλον ἑχάγωνον ἰσόπλευρόν τε καὶ ἰσογώνιον
ἐγγράψαι.}

{\greekfont Ἔστω ὁ δοθεὶς κύκλος ὁ ΑΒΓΔΕΖ; δεῖ δὴ εἰς τὸν ΑΒΓΔΕΖ
κύκλον ἑχάγωνον ἰσόπλευρόν τε καὶ ἰσογώνιον ἐγγράψαι.}

{\greekfont Ἤθθω τοῦ ΑΒΓΔΕΖ κύκλου διάμετρος ἡ ΑΔ, καὶ εἰλήφθω τὸ
κέντρον τοῦ κύκλου τὸ Η, καὶ κέντρῳ μὲν τῷ Δ διαστήματι δὲ
τῷ ΔΗ κύκλος γεγράφθω ὁ ΕΗΓΘ, καὶ ἐπιζευθθεῖσαι αἱ ΕΗ, ΓΗ
διήθθωσαν ἐπὶ τὰ Β, Ζ σημεῖα, καὶ ἐπεζεύθθωσαν αἱ ΑΒ, ΒΓ, ΓΔ,
ΔΕ, ΕΖ, ΖΑ; λέγω, ὅτι τὸ ΑΒΓΔΕΖ ἑχάγωνον ἰσόπλευρόν τέ
ἐστι καὶ ἰσογώνιον.}\\

\epsfysize=2.4in
\centerline{\epsffile{Book04/fig15g.eps}}

{\greekfont Ἐπεὶ γὰρ τὸ Η σημεῖον κέντρον ἐστὶ τοῦ ΑΒΓΔΕΖ κύκλου,
ἴση ἐστὶν ἡ ΗΕ τῇ ΗΔ. πάλιν, ἐπεὶ τὸ Δ σημεῖον κέντρον
ἐστὶ τοῦ ΗΓΘ κύκλου, ἴση ἐστὶν ἡ ΔΕ τῇ ΔΗ. ἀλλ
ἡ ΗΕ τῇ ΗΔ ἐδείθθη ἴση; καὶ ἡ ΗΕ ἄρα τῇ ΕΔ ἴση ἐστίν;
ἰσόπλευρον ἄρα ἐστὶ τὸ ΕΗΔ τρίγωνον; καὶ αἱ τρεῖς ἄρα α ὐτοῦ
γωνίαι αἱ ὑπὸ ΕΗΔ, ΗΔΕ, ΔΕΗ ἴσαι ἀλλήλαις εἰσίν,
ἐπειδήπερ τῶν ἰσοσκελῶν τριγώνων αἱ πρὸς τῇ βάσει γωνίαι
ἴσαι ἀλλήλαις εἰσίν; καί εἰσιν αἱ τρεῖς τοῦ τριγώνου γωνίαι
δυσὶν ὀρθαῖς ἴσαι; ἡ ἄρα ὑπὸ ΕΗΔ γωνία τρίτον ἐστὶ
δύο ὀρθῶν. ὁμοίως δὴ δειθθήσεται καὶ ἡ ὑπὸ ΔΗΓ τρίτον
δύο ὀρθῶν.
καὶ ἐπεὶ ἡ ΓΗ εὐθεῖα ἐπὶ τὴν ΕΒ σταθεῖσα
τὰς ἐφεχῆς γωνίας τὰς ὑπὸ ΕΗΓ, ΓΗΒ δυσὶν ὀρθαῖς
ἴσας ποιεῖ, καὶ λοιπὴ ἄρα ἡ ὑπὸ ΓΗΒ τρίτον ἐστὶ δύο
ὀρθῶν; αἱ ἄρα ὑπὸ ΕΗΔ, ΔΗΓ, ΓΗΒ γωνίαι ἴσαι
ἀλλήλαις εἰσίν; ὥστε καὶ αἱ κατὰ κορυφὴν α ὐταῖς αἱ ὑπὸ
ΒΗΑ, ΑΗΖ, ΖΗΕ ἴσαι εἰσὶν [ταῖς ὑπὸ ΕΗΔ, ΔΗΓ, ΓΗΒ].
αἱ ἓχ ἄρα γωνίαι αἱ ὑπὸ ΕΗΔ, ΔΗΓ, ΓΗΒ, ΒΗΑ, ΑΗΖ, ΖΗΕ
ἴσαι ἀλλήλαις εἰσίν. αἱ δὲ ἴσαι γωνίαι ἐπὶ ἴσων περιφερειῶν
βεβήκασιν; αἱ ἓχ ἄρα περιφέρειαι αἱ ΑΒ, ΒΓ, ΓΔ, ΔΕ, ΕΖ, ΖΑ
ἴσαι ἀλλήλαις εἰσίν. ὑπὸ δὲ τὰς ἴσας περιφερείας αἱ ἴσαι
εὐθεῖαι ὑποτείνουσιν; αἱ ἓχ ἄρα εὐθεῖαι ἴσαι 
ἀλλήλαις εἰσίν; ἰσόπλευρον ἄρα ἐστὶ το ΑΒΓΔΕΖ ἑχάγωνον.
λέγω δή, ὅτι καὶ ἰσογώνιον. ἐπεὶ γὰρ ἴση ἐστὶν ἡ ΖΑ
περιφέρεια τῇ ΕΔ περιφερείᾳ, κοινὴ προσκείσθω ἡ ΑΒΓΔ περιφέρεια; ὅλη ἄρα ἡ ΖΑΒΓΔ ὅλῃ τῇ ΕΔΓΒΑ  ἐστιν ἴση; καὶ βέβηκεν
ἐπὶ μὲν τῆς ΖΑΒΓΔ περιφερείας ἡ ὑπὸ 
ΖΕΔ γωνία, ἐπὶ
δὲ τῆς ΕΔΓΒΑ περιφερείας ἡ ὑπὸ ΑΖΕ γωνία;  ἴση 
ἄρα ἡ ὑπὸ 
ΑΖΕ γωνία τῇ ὑπὸ ΔΕΖ.  ὁμοίως δὴ δειθθήσεται, ὅτι 
καὶ αἱ λοιπαὶ γωνίαι τοῦ ΑΒΓΔΕΖ ἑχαγώνου κατὰ μίαν 
ἴσαι εἰσὶν ἑκατέρᾳ τῶν ὑπὸ ΑΖΕ, ΖΕΔ γωνιῶν; ἰσογώνιον
ἄρα ἐστὶ τὸ ΑΒΓΔΕΖ ἑχάγωνον. ἐδείθθη δὲ καὶ ἰσόπλευρον;
καὶ ἐγγέγραπται εἰς τὸν ΑΒΓΔΕΖ κύκλον.}

{\greekfont Εἰς ἄρα τὸν δοθέντα κύκλον ἑχάγωνον ἰσόπλευρόν τε καὶ
ἰσογώνιον ἐγγέγραπται; ὅπερ ἔδει ποιῆσαι.}\\~\\~\\~\\~\\

\begin{center}
{\large {\greekfont Πόρισμα}.}
\end{center}\vspace*{-7pt}

{\greekfont Ἐκ δὴ τούτου φανερόν, ὅτι ἡ τοῦ ἑχαγώνου πλευρὰ ἴση ἐστὶ
τῇ ἐκ τοῦ κέντρου τοῦ κύκλου.}

{\greekfont Ὁμοίως δὲ τοῖς ἐπὶ τοῦ πενταγώνου ἐὰν διὰ τῶν κατὰ
τὸν κύκλον διαιρέσεων ἐφαπτομένας τοῦ κύκλου ἀγάγωμεν,
περιγραφήσεται περὶ τὸν κύκλον ἑχάγωνον ἰσόπλευρόν τε
καὶ ἰσογώνιον ἀκολούθως τοῖς ἐπὶ τοῦ πενταγώνου
εἰρημένοις. καὶ ἔτι διὰ τῶν ὁμοίων τοῖς
ἐπὶ τοῦ πενταγώνου εἰρημένοις εἰς τὸ δοθὲν ἑχάγωνον κύκλον
ἐγγράψομέν τε καὶ περιγράψομεν; ὅπερ ἔδει ποιῆσαι.}}

\ParallelRText{
\begin{center}
{\large Proposition 15}
\end{center}

To inscribe an equilateral and equiangular hexagon
in a given circle.

Let $ABCDEF$ be the given circle. So it is required to inscribe an equilateral
and equiangular hexagon in circle $ABCDEF$.

Let the diameter $AD$ of circle $ABCDEF$ be drawn,$^\dag$ and let the center $G$ of
the circle be found [Prop.~3.1].
And let the circle $EGCH$ be drawn, with center $D$, and radius
$DG$.  And $EG$ and $CG$ being joined, let them be drawn across (the
circle) to points $B$ and $F$ (respectively). 
And let $AB$, $BC$, $CD$, $DE$, $EF$, and $FA$ be joined. 
I say that the hexagon $ABCDEF$ is equilateral and equiangular.

\epsfysize=2.4in
\centerline{\epsffile{Book04/fig15e.eps}}

For since point $G$ is the center of circle $ABCDEF$, $GE$ is equal to $GD$.
Again, since point $D$ is the center of circle $GCH$, $DE$ is equal to
$DG$. But, $GE$ was shown (to be) equal to $GD$. Thus, $GE$ is also equal to
$ED$. Thus, triangle $EGD$ is equilateral. Thus, its three angles
$EGD$, $GDE$, and $DEG$ are also equal to one another, inasmuch as the angles at the base of 
 isosceles triangles are equal to one another [Prop.~1.5]. And the three angles of the
 triangle are equal to
 two right-angles [Prop.~1.32].
 Thus, angle $EGD$ is one third of two right-angles. So, similarly,
  $DGC$ can also be shown (to be) one third of two right-angles.
   And since the straight-line $CG$, standing on $EB$, makes 
 adjacent angles $EGC$ and $CGB$ equal to two right-angles [Prop.~1.13], the remaining angle
 $CGB$ is thus also one third of two right-angles.
 Thus, angles $EGD$, $DGC$, and $CGB$ are equal to one another.
 And hence the (angles) opposite to them $BGA$, $AGF$, and $FGE$
 are also equal [to $EGD$, $DGC$, and $CGB$ (respectively)] [Prop.~1.15]. Thus,
 the six angles $EGD$, $DGC$, $CGB$, $BGA$, $AGF$, and $FGE$ are equal to one
 another. And equal angles stand on equal circumferences [Prop.~3.26]. Thus, the six circumferences
 $AB$, $BC$, $CD$, $DE$, $EF$, and $FA$ are equal to one another. And
  equal circumferences
  are subtended by equal straight-lines [Prop.~3.29]. Thus, the six
 straight-lines ($AB$, $BC$, $CD$, $DE$, $EF$, and $FA$) are equal to one another.
 Thus, hexagon $ABCDEF$ is equilateral. 
 So, I say that (it is) also equiangular. For since circumference $FA$
 is equal to circumference $ED$, let circumference $ABCD$ be added
 to both. Thus, the whole of $FABCD$ is equal to the whole of
 $EDCBA$. And angle $FED$ stands on circumference $FABCD$,
 and angle $AFE$ on circumference $EDCBA$. Thus, angle
 $AFE$ is equal to $DEF$ [Prop.~3.27].
 Similarly, it can also be shown that the remaining angles of hexagon
 $ABCDEF$ are individually equal to each of the angles $AFE$ and $FED$.
 Thus, hexagon $ABCDEF$ is equiangular. And it was also shown (to be)
 equilateral. And it has been inscribed in circle $ABCDE$.
 
 Thus, an equilateral and equiangular hexagon has been inscribed
 in the given circle. (Which is) the very thing it was required to do.\\
 
 \begin{center}
{\large Corollary}
\end{center}\vspace*{-7pt}

So, from this, (it is) manifest that a side of the hexagon is equal
to the radius of the circle.

And similarly to  a pentagon, if we draw  tangents to the circle through
the (sixfold) divisions of the (circumference of the) circle, an equilateral and equiangular hexagon
can be circumscribed about the circle, analogously to the aforementioned pentagon. And, further, by (means) similar to the aforementioned pentagon,
we can inscribe and circumscribe a circle in (and about) a given hexagon. 
(Which
is) the very thing it was required to do.}
\end{Parallel}
{\footnotesize \noindent$^\dag$ See the
footnote to Prop.~4.6.} 

%%%%%%
% Prop 4.16
%%%%%%
\pdfbookmark[1]{Proposition 4.16}{pdf4.16}
\begin{Parallel}{}{} 
\ParallelLText{
\begin{center}
{\large \ggn{16}.}
\end{center}\vspace*{-7pt}

{\greekfont Εἰς τὸν δοθέντα κύκλον πεντεκαιδεκάγωνον ἰσόπλευρ-όν τε καὶ
ἰσογώνιον ἐγγράψαι.}

\epsfysize=2.2in
\centerline{\epsffile{Book04/fig16g.eps}}

{\greekfont Ἔστω ὁ δοθεὶς κύκλος ὁ ΑΒΓΔ; δεῖ δὴ εἰς τὸν ΑΒΓΔ
κύκλον πεντεκαιδεκάγωνον ἰσόπλευρόν τε καὶ ἰσογώνιον
ἐγγράψαι.}

{\greekfont Ἐγγεγράφθω εἰς τὸν ΑΒΓΔ κύκλον τριγώνου μὲν ἰσοπλεύρου
τοῦ εἰς α ὐτὸν ἐγγραφομένου πλευρὰ ἡ ΑΓ, πενταγώνου δὲ
ἰσοπλεύρου ἡ ΑΒ; οἵων ἄρα ἐστὶν ὁ ΑΒΓΔ κύκλος ἴσων
τμήματων δεκαπέντε, τοιούτων ἡ μὲν ΑΒΓ περιφέρεια τρίτον
οὖσα τοῦ κύκλου ἔσται πέντε, ἡ δὲ ΑΒ περιφέρεια πέμτον
οὖσα τοῦ κύκλου ἔσται τριῶν; λοιπὴ ἄρα ἡ ΒΓ τῶν
ἴσων δύο. τετμήσθω ἡ ΒΓ δίθα κατὰ τὸ Ε; ἑκατέρα ἄρα τῶν
ΒΕ, ΕΓ περιφερειῶν πεντεκαιδέκατόν ἐστι τοῦ ΑΒΓΔ κύκλου.}

{\greekfont Ἐὰν ἄρα ἐπιζεύχαντες τὰς ΒΕ, ΕΓ ἴσας α ὐταῖς κατὰ 
τὸ συνεθὲς εὐθείας ἐναρμόσωμεν εἰς τὸν ΑΒΓΔ[Ε]
κύκλον, ἔσται εἰς α ὐτὸν ἐγγεγραμμένον πεντεκαιδεκάγωνον
ἰσόπλευ-ρόν τε καὶ ἰσογώνιον; ὅπερ ἔδει ποιῆσαι.}

{\greekfont Ὁμοίως δὲ τοῖς ἐπὶ τοῦ πενταγώνου ἐὰν διὰ τῶν κατὰ
τὸν κύκλον διαιρέσεων  ἐφαπτομένας τοῦ  κύκλου ἀγάγωμεν,
περιγραφήσεται περὶ τὸν κύκλον πεντεκαιδεκάγωνον ἰσόπλευρόν
τε καὶ ἰσογώνιον. ἔτι δὲ διὰ τῶν ὁμοίων τοῖς
ἐπὶ τοῦ πενταγώνου δείχεων καὶ εἰς τὸ δοθὲν πεντεκαιδεκάγωνον
κύκλον ἐγγράψομέν τε καὶ περιγράψομεν; ὅπερ ἔδει ποιῆσαι.}}

\ParallelRText{
\begin{center}
{\large Proposition 16}
\end{center}

To inscribe an equilateral and equiangular fifteen-sided figure in a given circle.

\epsfysize=2.2in
\centerline{\epsffile{Book04/fig16e.eps}}

Let $ABCD$ be the given circle. So it is required to inscribe an equilateral
and equiangular fifteen-sided figure in circle $ABCD$.

Let the side $AC$ of an equilateral triangle inscribed in (the circle) [Prop.~4.2],
and (the side) $AB$ of an (inscribed) equilateral pentagon [Prop.~4.11],
 be inscribed in circle $ABCD$.
 Thus,   just as the circle $ABCD$ is (made up) of fifteen equal pieces, the
 circumference $ABC$, being a third of the circle, will be (made up) of five
 such (pieces), and the
 circumference $AB$, being a fifth of the circle, will be (made up) of  three.
 Thus, the remainder $BC$ (will be made up) of two equal (pieces).
 Let (circumference) $BC$ be cut in half at $E$ [Prop.~3.30].
 Thus, each of the circumferences $BE$ and $EC$ is one fifteenth of the circle
 $ABCDE$.
 
 Thus, if, joining $BE$ and $EC$, we continuously insert straight-lines
 equal to them into circle $ABCD[E]$ [Prop.~4.1],
 (then) an equilateral and equiangular fifteen-sided figure will be
 inserted into (the circle). (Which is) the very thing it was required to do.
 
 And similarly to the pentagon, if we draw  tangents to the circle through
the (fifteenfold) divisions of the (circumference of the) circle, we can
circumscribe an equilateral and equiangular fifteen-sided figure about
the circle. 
 And, further, through similar proofs to the  pentagon,
 we can also inscribe and circumscribe  a circle in (and about)  a given fifteen-sided 
  figure. (Which is) the very thing it was required to do.}
\end{Parallel}
